% -*- latex -*-
\documentclass[a4paper,12pt]{article}
\usepackage[utf8]{inputenc}
\usepackage[T1]{fontenc}
\usepackage[ngerman]{babel}
% \usepackage{lmodern}
\usepackage{amsmath,amssymb,amsthm}
\usepackage{geometry}
\usepackage{graphicx}
\usepackage[hidelinks]{hyperref}
\usepackage{booktabs}
\usepackage{listings}
\usepackage{xcolor}
\usepackage{enumitem}
\usepackage{float}
\usepackage{placeins}
\usepackage{array}
\newcolumntype{P}[1]{>{\raggedright\arraybackslash}p{#1}}
% \usepackage{microtype}
% Zeilenumbruch: verhindert Überläufe über den rechten Rand, wo die Silbentrennung nicht reicht.
\setlength{\emergencystretch}{12em}
\geometry{margin=2.5cm}
\setlength{\tabcolsep}{3pt}

\definecolor{mainblue}{RGB}{25, 70, 140}
\definecolor{accentred}{RGB}{180, 50, 30}
\definecolor{lightgray}{RGB}{245, 245, 248}
\definecolor{darkgray}{RGB}{60, 60, 70}
\definecolor{darkgreen}{RGB}{30, 100, 50}
 
\hypersetup{colorlinks=true, linkcolor=mainblue, citecolor=accentred, urlcolor=mainblue}

\title{\textbf{Wann ist das Datennetz gut genug für KI?}\\\large Eine ökonomische und formale Betrachtung des AI Data Network}
\author{\textsc{Stephan Epp}\\ {\small \texttt{\textbf{M.Sc.}}}}
\date{15. Oktober 2026}

% Theorem environments
\theoremstyle{definition}
\newtheorem{definition}{Definition}[section]
\newtheorem{bemerkung}[definition]{Bemerkung}
\newtheorem{satz}[definition]{Satz}
\newtheorem{lemma}[definition]{Lemma}
\newtheorem{korollar}[definition]{Korollar}

% Notation
\lstset{basicstyle=\ttfamily\footnotesize,breaklines=true,frame=single,backgroundcolor=\color{lightgray},columns=fullflexible,keepspaces=true,inputencoding=utf8,extendedchars=true,literate={ä}{{\"a}}1 {ö}{{\"o}}1 {ü}{{\"u}}1 {Ä}{{\"A}}1 {Ö}{{\"O}}1 {Ü}{{\"U}}1 {ß}{{\ss}}1}
\newcommand{\bits}{\{0,1\}}
\newcommand{\NN}{\mathbb{N}}
\newcommand{\RR}{\mathbb{R}}
\newcommand{\QQ}{\mathbb{Q}}
\newcommand{\cP}{\mathrm{P}}
\newcommand{\cNP}{\mathrm{NP}}
\newcommand{\E}{\mathbb{E}}
\newcommand{\Var}{\mathrm{Var}}
\newcommand{\KL}{\mathrm{KL}}

\begin{document}
\maketitle

\begin{abstract}
Diese Arbeit untersucht, wann das Datennetz eines Unternehmens für eine konkrete KI-Anwendung gut genug ist. Ausgangspunkt ist nicht die Forderung nach fehlerfreien Daten, sondern eine ökonomische Abwägung: Ist der verbleibende Fehler im Verhältnis zu den Folgen der unterstützten Entscheidung vertretbar? Eine Fallstudie zur sinkenden Deckungsbeitragsmarge zeigt, wie eine plausibel klingende, aber unvollständige KI-Erklärung zu einer falschen Maßnahme führen kann, wenn Logistikdaten und Produktmix im Datennetz fehlen. Darauf aufbauend verbindet die Arbeit ein formales Modell von Datennetz, Data Product, Provenance und Error Budget mit einem praktischen Vorgehen. Sie beschreibt KI als Hinweis- und Prüfwerkzeug, vier stufenweise Vertrauensniveaus, ein Data Product als wiederverwendbare Grundlage, drei benannte Verantwortlichkeiten und messbare Kriterien für Fortschritt. Die zentrale These lautet: Abweichungen zwischen KI-Antwort und Unternehmensanalyse sind nicht nur ein Risiko, sondern können Datenlücken und zugleich mögliche Lücken in der eigenen Analyse sichtbar machen.
Um das Prüfverfahren in Unternehmensabläufe einzubinden, wird es zusätzlich als zustandslose JSON-Schnittstelle (FastAPI) bereitgestellt. Sie bildet 28 Operationen ab, liefert für gleiche Eingaben identische Ergebnisse wie die Bibliothek, behandelt Unbekanntes an der Schnittstelle nie als Erfolg, begrenzt Eingabegrößen gegen Überlastung und ist mit 291 Tests vollständig abgedeckt. Der Ausblick skizziert fünf Richtungen: empirische Validierung, Modellerweiterung, Ausbau der Schnittstelle, organisatorische Einbettung und die Anbindung von KI-Agenten.
\end{abstract}

\tableofcontents
\newpage

\section{Einleitung}
Die drei leitenden Gedanken dieser Arbeit sind eng miteinander verbunden. Erstens hängt Vertrauen in eine Antwort nicht davon ab, ob sie vollkommen fehlerfrei ist, sondern davon, ob ein möglicher Fehler für den konkreten Verwendungszweck klein genug ist. Zweitens bezeichnet das Datennetz mehr als die Summe gespeicherter Daten: Entscheidend sind auch Verknüpfungen, Zugriffswege, gemeinsame Definitionen und Aktualität. Dasselbe KI-Modell kann deshalb auf einem gut verbundenen Netz nützliche und auf einem fragmentierten Netz irreführende Antworten geben. Drittens können KI-Antworten bereits vor vollständiger Datenreife als Hinweise und Gegenprüfung der eigenen wirtschaftlichen Analyse nützlich sein.

Diese Perspektive vermeidet zwei Extreme: KI-Antworten weder blind als Wahrheit zu übernehmen noch den Einsatz bis zu einer unerreichbaren Perfektion der Daten aufzuschieben. Stattdessen werden Nutzung und Autonomie nach den möglichen Folgen eines Fehlers gestuft. Die Arbeit führt zunächst in die Lücke zwischen KI-Ambition und Datenreife ein und erläutert eine Fallstudie zur sinkenden Marge. Danach folgen ein formales Modell, eine ökonomische Bedingung für ausreichende Qualität und ein Lernkreislauf für Abweichungen. Abschließend werden Data Products, Verantwortlichkeiten, Umsetzungsschritte, Unternehmensgrößen, typische Stolperfallen, Messgrößen, ein deterministisches Prüfverfahren mit Freigabebericht und Regressionsvergleich sowie eine Go-live-Checkliste behandelt. Ein eigenes Kapitel beschreibt die Bereitstellung des Prüfverfahrens als JSON-Schnittstelle (Abschnitt~\ref{sec:api}); der Ausblick (Abschnitt~\ref{sec:outlook}) ordnet offene Aufgaben und Erweiterungen.

\section{Ausgangslage: KI-Ambition und Datenreife}
\label{sec:starting-point}

In der dem Konzept zugrunde liegenden Webinar-Präsentation werden drei Hindernisse für die Nutzung von KI hervorgehoben: eingeschränkter Datenzugriff, unzureichende Datenqualität und voneinander getrennte Systeme. Die dort genannten Umfragewerte dienen hier als Kontext und werden als Angaben aus dem Webinar wiedergegeben \cite{webinar2026}; sie sind keine in dieser Arbeit unabhängig überprüfte Primärerhebung.

\begin{table}[htbp]
\centering
\small
\caption{Im Webinar berichtete Angaben zu KI und Datenreife}
\label{tab:readiness}
\begin{tabular}{p{0.33\linewidth}p{0.57\linewidth}}
\toprule
\textbf{Aspekt} & \textbf{Im Webinar genannte Angabe} \\
\midrule
KI-Initiativen & 97\,\% der Unternehmen haben KI-Initiativen. \\
Vollständige Datenreife & 5\,\% halten ihre Daten für vollständig KI-bereit. \\
Hindernisse & Eingeschränkter Datenzugriff: 50\,\%; Datenqualität: 40\%; getrennte Systeme: 38\,\%. \\
KI-Agenten & 11\,\% setzen Agenten bereits ein; weitere 29\,\% planen den Einsatz. \\
Investitionen & Erfolgreiche Unternehmen investieren laut Präsentation bis zu viermal mehr in Datenqualität und Data Governance. \\
\bottomrule
\end{tabular}
\end{table}
\noindent\footnotesize Im Webinar genannte Quellen: Dun \& Bradstreet (10.000 Unternehmen in 32 Ländern, 2026), Bitkom (2026) und Gartner (2026). Die Zuordnung der Einzelwerte zu diesen Quellen wird in der Präsentation nicht näher aufgeschlüsselt.\normalsize

Diese Hindernisse sind nicht primär Modellprobleme. Eingeschränkter Zugriff bedeutet, dass relevante Daten zwar existieren, aber für die fragende Person oder die KI nicht erreichbar sind. Bei Qualitätsproblemen sind Daten vorhanden, aber beispielsweise unvollständig, verspätet oder nicht verlässlich genug. Getrennte Systeme verlieren Beziehungen zwischen Fakten: So kann ein Expressfrachtaufwand im Logistiksystem verbucht sein, ohne mit dem Auftrag verbunden zu sein, der ihn verursacht hat. Ein besseres Sprachmodell schließt diese Lücke nicht von selbst.

Damit lautet die passende Diagnose nicht nur „Ist das KI-Modell gut genug?“, sondern vor allem: „Ist unser Datennetz für diese konkrete Frage gut genug?“ Die zeitliche Lücke zwischen bereits laufenden oder geplanten Agenten und geringer eingeschätzter Datenreife macht eine solche Prüfung besonders wichtig.

\section{Fallstudie: Warum sinkt die Marge?}
\label{sec:case-study}

Die folgende illustrative Fallstudie greift die im Webinar präsentierte Situation eines Maschinenbauzulieferers für das erste bis dritte Quartal 2026 auf. Sie veranschaulicht, wie eine sachlich teilweise richtige KI-Antwort durch ein unvollständiges Datennetz zu einer unverhältnismäßigen Empfehlung führen kann.

\subsection{Ausgangslage und KI-Antwort}
Der Umsatz stieg gegenüber dem Vorjahr von 10{,}0\,Mio.\,€ auf 11{,}0\,Mio.\,€. Gleichzeitig sank die Deckungsbeitragsmarge II (DB\,II) von 28\,\% auf 24\,\%. Bei Anwendung der Vorjahresmarge auf den aktuellen Umsatz wären 3{,}08\,Mio.\,€ DB\,II erzielt worden; tatsächlich waren es 2{,}64\,Mio.\,€. Die Differenz beträgt 440\,T€.

Auf die Frage „Warum sinkt unsere Marge, und was sollten wir tun?“ verwies die KI vor allem auf steigende Rabatte. Die durchschnittliche Rabattquote sei von 6{,}1\,\% auf 6{,}7\,\% gestiegen, besonders bei Großkunden. Als Maßnahmen schlug sie Genehmigungen für Rabatte über 5\,\%, Neuverhandlungen mit den zehn größten Kunden und eine monatliche Berichterstattung vor. Die Antwort klang plausibel und konkret, enthielt aber weder nachvollziehbare Quellenangaben noch einen Datenstichtag.

\subsection{Ergebnis der Controlling-Analyse}
Die nachfolgende interne Analyse ordnete den Margenabstand drei Ursachen zu. Der Expressfrachtfall betraf 31 teure Eillieferungen, die nach einem Engpass bei einem Lieferantenteil erforderlich waren, um den Kunden Nord termingerecht zu beliefern.

\begin{table}[htbp]
\centering
\small
\caption{Aufteilung der DB-II-Abweichung in der Fallstudie}
\label{tab:margin-case}
\begin{tabular}{p{0.24\linewidth}p{0.10\linewidth}p{0.14\linewidth}p{0.16\linewidth}p{0.25\linewidth}}
\toprule
\textbf{Ursache} & \textbf{Betrag} & \textbf{Anteil} & \textbf{Von KI erkannt?} & \textbf{Grund} \\
\midrule
Expressfracht & 230\,T€ & ca. 52\,\% & Nein & Im Logistiksystem gebucht, für die KI nicht sichtbar oder nicht mit Aufträgen verbunden. \\
Produktmix & 150\,T€ & ca. 34\,\% & Nein & Die Produktgruppe fehlte im zugänglichen Datennetz. \\
Rabatte & 60\,T€ & ca. 14\,\% & Ja & In den zugänglichen Vertriebsdaten sichtbar. \\
\midrule
\textbf{Gesamt} & \textbf{440\,T€} & \textbf{100\,\%} & & \\
\bottomrule
\end{tabular}
\end{table}

Kunde Nord erwirtschaftete 2{,}0\,Mio.\,€ Umsatz. Die teuren Eillieferungen waren nicht durch ein Fehlverhalten des Kunden verursacht worden. Eine auf die KI-Empfehlung gestützte Rabattverhandlung hätte daher einen wertvollen Kunden unter Druck setzen können, während die größte Ursache, die Expressfracht, unangetastet geblieben wäre.

\subsection{Lehren aus dem Fall}
Die KI-Antwort war nicht vollständig falsch: Der Rabattanteil war tatsächlich gestiegen. Sie erklärte jedoch nur rund 14\,\% der untersuchten Differenz und ließ die beiden größeren Ursachen außer Acht. Die Fallstudie macht fünf Punkte deutlich:
\begin{enumerate}
\item Eine flüssige und plausible Antwort ist nicht gleichbedeutend mit einer vollständigen Erklärung.
\item Die zentralen blinden Flecken lagen im Datennetz: fehlende Verknüpfung von Fracht und Auftrag sowie eine nicht verfügbare Produktgruppendimension. Ein anderes Modell hätte dieselben Datenlücken vorgefunden.
\item Ohne Definition, Quelle und Zeitstempel lässt sich die Antwort kaum angemessen einordnen.
\item Eine ungeprüfte Empfehlung kann neben dem unmittelbaren Zahlenfehler auch Geschäftsbeziehungen und künftige Umsätze gefährden.
\item Als Hinweis war die Beobachtung zur Rabattentwicklung dennoch nützlich. Gefährlich wäre erst, sie als abschließende Gesamterklärung zu behandeln.
\end{enumerate}
Für eine erste, informelle Auffälligkeit kann ein erklärter Anteil von 14\,\% ein Anlass zur weiteren Prüfung sein. Für eine Handlungsempfehlung, die einen Kunden mit 2{,}0\,Mio.\,€ Umsatz betrifft, ist die Erklärung offenkundig unzureichend. Das Netz war somit nicht allgemein „schlecht“, sondern für diese konkrete Frage und diesen Verwendungszweck nicht ausreichend.

\section{Formales Modell des Datennetzes}\label{sec:formal}
Wir modellieren die Unternehmensdatenlandschaft als gerichteten multigraphischen Datenraum und definieren Data Products als wohldefinierte Subnetze mit Metadaten.

\begin{definition}[Daten-Netzwerk]
Sei $G=(V,E)$ ein gerichteter Multigraph, wobei Knoten $v\in V$ Datentische, Datensätze oder Dokumente repräsentieren und Kanten $e=(u\to v,\varphi)$ eine join-/mapping-Relation mit einer Abbildungsfunktion/Transformation $\varphi$ kennzeichnen. Jeder Knoten trägt Attribute (Spalten) $A_v$ und eine Qualität $q_v\in[0,1]$ sowie eine Aktualität $t_v\in\RR_{\ge0}$ (letzte Aktualisierung).
\end{definition}

\begin{definition}[Data Product]
Ein \emph{Data Product} $P$ ist ein Tupel $(V_P,E_P,\mathcal M,\ell)$ mit Teilmenge $V_P\subseteq V$, $E_P\subseteq E$, einer semantischen Beschreibung $\mathcal M$ (Meaning: Definitionen, Synonyme, Regeln) und einem Label $\ell$ (Owner, Quality, Frequency, Version). $P$ ist \emph{geschlossen} in dem Sinne, dass alle Kanten für die relevanten Aggregationen innerhalb $P$ definiert sind oder durch referenzielle Metadaten nachweisbar verbunden werden können.
\end{definition}

Die KI-Anwendung ist ein Algorithmus $\mathcal A$ der auf Queries $q$ und Zugang zu einem (Teil-)Netz $G'$ eine Antwort $\hat y=\mathcal A(q;G')$ liefert. Die Wahrheit (ground truth) zu einer Frage ist eine Funktion $y^*=f_q(G_{true})$ über das (hypothetisch) vollständige Datennetz $G_{true}$. In der Praxis ist $G'\subseteq G_{true}$.

\begin{definition}[Provenance]
Zur Antwort $\hat y$ liefert das System eine Provenance-Menge $\Pi(\hat y)$ bestehend aus: (i) referenzierten Data Products/Tabellen, (ii) Versionen und Zeitstempel, (iii) definitorischen Regeln aus $\mathcal M$ die angewendet wurden.
\end{definition}

\subsection{Einflussgrößen und beobachtbare Variablen}
Für eine Frage $q$ definiert ein Domänenexperte die Menge der \emph{wesentlichen Einflussgrößen} $S_q=\{s_1,\dots,s_k\}$, die die Antwort beeinflussen (z. B. Discount, Freight, ProductGroup). Jede $s_i$ ist eine Funktion auf Knoten/Attributen des Netzes: $s_i=\phi_i(G_{true})$. In $G'$ sind nur Submengen $\tilde S_q$ beobachtbar.

\section{Fehlerzusammensetzung und Gut genug}
\label{sec:error}
Wir wollen für eine konkrete Entscheidung die zulässige Fehlergrösse analysieren.

\begin{definition}[Entscheidungs-Stakes und Error Budget]
Sei $D$ eine Entscheidung mit monetärem Stake $W\ge0$ (z. B. erwarteter Einfluss auf Marge). Ein erwarteter Fehler in der Antwort $\hat y$ führt zu einem erwarteten monetären Schaden $L(\hat y,y^*)$. Ein Error Budget ist $\varepsilon>0$ so dass $\E[L(\hat y,y^*)]\le\varepsilon W$ nötig ist, damit die KI-Antwort für die automatisierte oder vorbereitende Nutzung \emph{zulässig} ist.
\end{definition}

Typische Loss-Funktionen sind absolutes Delta in Euro, Rangdifferenz, oder Entscheidungsverlust (z. B. falsche Top-K-Kunden).

\begin{satz}[Fehlerzerlegung]
Sei $q$ eine Frage mit wahren Treibern $S_q$ und beobachteter Menge $\tilde S_q\subseteq S_q$ im genutzten Datennetz $G'$. Unter milden Regularitätsannahmen (lineare Einflussbeziehung, additive Störgröße) zerfällt der Antwortfehler in:
\[
\hat y - y^* = B_{miss} + B_{def} + \eta
\]
wobei
\begin{itemize}
\item $B_{miss}$ der Bias wegen \emph{fehlender Einflussgrößen} ($S_q\setminus\tilde S_q$) ist,
\item $B_{def}$ der Bias aus \emph{uneinheitlicher Definition} (Inkompatible Metriken) ist,
\item $\eta$ ein zufälliger Fehler (Rauschen) aus Datenqualität/Update-Verzögerung.
\end{itemize}
Der Erwartungswert $\E[\hat y - y^*] = B_{miss}+B_{def}$.
\end{satz}
\begin{proof}
Unter der Annahme, dass $y^*=\sum_{i\in S_q} \alpha_i s_i + \xi$ (lineares Modell) und $\hat y=\sum_{i\in\tilde S_q}\alpha_i s_i + \tilde\xi$ (Modell auf beobachteten Variablen), folgt:
\[
\hat y-y^*= -\sum_{i\in S_q\setminus\tilde S_q} \alpha_i s_i + (\tilde\xi-\xi) + \Delta_{def}
\]
Setze $B_{miss}=-\sum_{i\in S_q\setminus\tilde S_q}\alpha_i s_i$, $\eta=\tilde\xi-\xi$ und $B_{def}=\Delta_{def}$; das ergibt die Behauptung.
\end{proof}

Damit ist die Kernaufgabe: Für gegebene Stakes $W$ schätzen, wie groß $|B_{miss}|+|B_{def}|+\mathrm{std}(\eta)$ werden darf.

\subsection{Gut genug als ökonomische Bedingung}
\begin{satz}[Ökonomische Gut-Genug-Bedingung]
Die KI-Antwort ist für eine Entscheidung mit Stakes $W$ "gut genug" genau dann, wenn die erwartete monetäre Auswirkung des Fehlers kleiner ist als das Error Budget, d. h.
\[
\E[L(\hat y,y^*)] \le \varepsilon W.
\]
Insbesondere gilt unter linearem Loss $L(\hat y,y^*)=c|\hat y-y^*|$ die praktische Bedingung
\[
c\,\E[|B_{miss}+B_{def}+\eta|] \le \varepsilon W.
\]
\end{satz}
\begin{proof}
Direkt aus Definition des Loss und Error Budget.
\end{proof}

Diese Bedingung ist operationalisierbar: Man misst erklärten Anteil (explained share) auf Testfragen mit bekannter Referenz und vergleicht mit akzeptiertem Anteil.

\subsection{Vertrauensstufen und Fehlerbudget}
Ein Error Budget macht die zulässige Abweichung für einen Anwendungsfall explizit. Es ist keine globale Qualitätsnote für ein Data Product: Dieselben Daten können für eine explorative Trendfrage genügen und für einen Eingriff in einen wichtigen Kundenvertrag unzureichend sein. Entscheidend sind der Zweck, die mögliche Schadenshöhe, die menschliche Aufsicht und die Messbarkeit des Fehlers.

\begin{table}[htbp]
\centering
\small
\caption{Nutzungsstufen nach erforderlicher Zuverlässigkeit}
\label{tab:reliability-levels}
\begin{tabular}{p{0.08\linewidth}p{0.25\linewidth}p{0.37\linewidth}p{0.21\linewidth}}
\toprule
\textbf{Stufe} & \textbf{Rolle der KI} & \textbf{Anforderung an Datennetz und Kontrolle} & \textbf{Beispiel} \\
\midrule
1 & Hinweis & Geringere Abdeckung ist vertretbar; der Mensch prüft die Spur eigenständig. & Rabattentwicklung genauer untersuchen. \\
2 & Gegenprüfung & Zentrale Kennzahlen sind verknüpft und definiert; Abweichungen werden verglichen. & KI-Aufschlüsselung mit Controlling-Analyse abgleichen. \\
3 & Vorbereitung von Entscheidungen & Qualität ist messbar; Quelle, Definition und Zeitstempel sind verfügbar; ein Mensch entscheidet. & Begründete Empfehlung für die Geschäftsleitung. \\
4 & Automatisierte Handlung & Sehr hohe Zuverlässigkeit, ausdrückliche Freigabe, Protokollierung und klarer Eingriffspfad. & KI legt nach Freigabe eine Aufgabe an. \\
\bottomrule
\end{tabular}
\end{table}

Organisationen müssen nicht mit der höchsten Stufe beginnen. Hinweise und Gegenprüfungen können früh Nutzen bringen, während Testabdeckung, Datenverknüpfungen und Governance schrittweise verbessert werden. Eine höhere Stufe setzt einen belegbaren Reifegrad voraus; sie folgt nicht automatisch daraus, dass die vorherige Stufe eine Zeit lang funktioniert hat.

\subsection{Stake-bezogene Anforderungen}
Die nachfolgende Einordnung ist eine Entscheidungshilfe, keine allgemeingültige numerische Norm. Die verantwortliche Person muss Fehlermaß, Referenz und akzeptable Schwelle für jeden konkreten Anwendungsfall festlegen.

\begin{table}[htbp]
\centering
\small
\caption{Beispiele für Stakes und Anforderungen an das Datennetz}
\label{tab:stakes}
\begin{tabular}{p{0.34\linewidth}p{0.20\linewidth}p{0.19\linewidth}p{0.19\linewidth}}
\toprule
\textbf{Fragestellung} & \textbf{Typische Stakes} & \textbf{Zulässiger Fehler} & \textbf{Erforderliche Grundlage} \\
\midrule
Welche Produktgruppen verkaufen sich am besten? & Niedrig, informativ & Eher großzügig & Kerndaten genügen; Quelle und Zeitraum bleiben nachvollziehbar. \\
Wo verlieren wir Marge und warum? & Mittel bis hoch & Haupttreiber müssen erfasst sein & Wesentliche Einflussgrößen müssen definiert und verbunden sein. \\
Mit welchen Kunden verhandeln wir neu? & Hoch, betrifft Beziehungen und Umsatz & Gering & Kausale Gesamtsicht und Qualitätsnachweis erforderlich. \\
Automatisch eine Aufgabe oder Maßnahme auslösen & Sehr hoch, gegebenenfalls ohne unmittelbare Prüfung & Sehr gering & Governance, Genehmigung, Protokollierung und Eskalationsweg. \\
\bottomrule
\end{tabular}
\end{table}

Für den Fall aus Abschnitt~\ref{sec:case-study} ist daher zwischen einer explorativen Frage und der Empfehlung „Was sollen wir tun?“ zu unterscheiden. Die erste kann eine unvollständige Antwort als Suchhinweis verwenden; die zweite verlangt Abdeckung der entscheidenden Treiber und eine Prüfung möglicher Nebenwirkungen. Ein einfacher Maßstab ist, den auf Testfällen belegten Anteil des Effekts der konkreten Entscheidung gegenüberzustellen und zu fragen, ob die nicht erklärte Restgröße noch zu groß ist.

\section{Abweichungen als Chance}\label{sec:opportunity}
Die besondere These dieser Arbeit: Eine Abweichung zwischen AI-Antwort und interner Analyse kann ein Indikator sein, der Information liefert und damit eine \emph{Chance} darstellt.

\subsection{KI als Hinweis und Gegenprüfung}
Ein Hinweis ist kein Urteil. Der Nutzen einer KI-Antwort kann darin liegen, eine bislang übersehene Richtung oder Auffälligkeit kenntlich zu machen, die anschließend fachlich geprüft wird. Im Margenbeispiel war der Anstieg der Rabattquote ein wahrer und prüfenswerter Hinweis, obwohl er nur einen kleinen Teil des gesamten Abstands erklärte.

Als Gegenprüfung wird die KI-Antwort mit der eigenen Strategie, Erfahrung oder Controlling-Analyse verglichen. Stimmen die Ergebnisse überein, kann das die Sicherheit in beiden Analysen erhöhen. Weichen sie voneinander ab, ist die Differenz ein Anlass, die Ursachen in Daten, Verknüpfungen und Definitionen zu untersuchen. Der Vergleich kann zugleich zeigen, dass die KI zusätzliche Aspekte gefunden hat, die in der bisherigen Unternehmensanalyse fehlten.

\subsection{Formalisierung als diagnostisches Signal}
Wir modellieren die Situation probabilistisch. Sei $H_0$ die Hypothese "Kein fehlender Treiber" (\(S_q = \tilde S_q\) und Definitionsgleichheit) und $H_1$ die Alternative "Es fehlt mindestens ein Treiber oder Definitionskonflikt".

Die KI erzeugt eine Antwort $\hat y$. Der interne Controlling-Wert (Referenz) ist $y_c$ (menschliche Analyse). Beobachte die Abweichung $d=\hat y-y_c$.

\begin{satz}[Abweichung als Bayesscher Indikator]
Angenommen es existieren Likelihoods $p(d\mid H_0)$ und $p(d\mid H_1)$. Dann ist die Bayes-Posterior Odds
\[
\frac{P(H_1\mid d)}{P(H_0\mid d)} = \frac{P(H_1)}{P(H_0)}\cdot\frac{p(d\mid H_1)}{p(d\mid H_0)}.
\]
Ist $p(d\mid H_1)>p(d\mid H_0)$, so erhöht die Beobachtung $d$ die posteriori Wahrscheinlichkeit von $H_1$ gegenüber dem Prior. Im Erwartungswert kann eine systematische Auswertung solcher Abweichungen die Unsicherheit über den Zustand des Datennetzes verringern; ob daraus ein wirtschaftlicher Nutzen entsteht, hängt zusätzlich von den möglichen Maßnahmen und ihren Kosten ab.
\end{satz}
\begin{proof}
Bayessches Theorem.
\end{proof}

\begin{bemerkung}
In Worten: Wenn eine bestimmte Abweichung unter \(H_0\) unwahrscheinlicher als unter \(H_1\) ist, verschiebt sie die Einschätzung in Richtung fehlender Treiber oder Definitionskonflikte. Das ist ein diagnostischer Hinweis, aber kein Beweis: Auch eine fehlerhafte eigene Analyse, ein abweichender Datenstand oder eine andere zulässige Abgrenzung kann die Differenz erklären.
\end{bemerkung}

\subsection{Abweichungen systematisch diagnostizieren}
Die Abweichung sollte nicht vorschnell dem Modell angelastet werden. Stattdessen werden Ursachen geprüft und jeweils mit einer konkreten Reaktion verknüpft:
\begin{table}[htbp]
\centering
\small
\caption{Mögliche Ursachen einer Abweichung und geeignete Reaktionen}
\label{tab:deviation-causes}
\begin{tabular}{p{0.21\linewidth}p{0.31\linewidth}p{0.40\linewidth}}
\toprule
\textbf{Mögliche Ursache} & \textbf{Typisches Anzeichen} & \textbf{Reaktion} \\
\midrule
Fehlende Datenquelle & Ein Treiber wie Fracht ist in einem anderen System vorhanden. & Quelle mit dem Data Product verbinden und Zuständigkeit festlegen. \\
Fehlendes Attribut & Eine wichtige Dimension wie Produktgruppe ist nicht verfügbar. & Attribut ergänzen und dessen Vollständigkeit überwachen. \\
Unklare Definition & Beteiligte verstehen unter „Marge“ unterschiedliche Kennzahlen. & Data Owner legt Bedeutung fest und dokumentiert Synonyme und Geltungsbereich. \\
Veraltete Daten & Antwort passt zu einem älteren Zustand oder Stichtag. & Aktualisierungsfrequenz klären und Datenstichtag ausweisen. \\
Unvollständige eigene Analyse & KI macht einen zuvor übersehenen Treiber sichtbar. & Hinweis prüfen und gegebenenfalls die interne Analyse korrigieren. \\
\bottomrule
\end{tabular}
\end{table}
Die Gegenprüfung funktioniert somit in beide Richtungen: Sie kann Mängel im Datennetz aufdecken und zugleich die Annahmen oder Berechnungen des Unternehmens in Frage stellen.

\subsection{Konkretes Lern-Loop-Modell}
Sei $\theta$ ein latenter Zustand des Datennetzes (zum Beispiel, welche Quellen verknüpft sind). Vor der Beobachtung hat das Unternehmen einen Prior $\pi_0(\theta)$. Eine Abweichung $d$ liefert die Likelihood $p(d\mid\theta)$ und damit den Posterior $\pi_1(\theta)\propto p(d\mid\theta)\pi_0(\theta)$. Die Wahrscheinlichkeit von Zuständen mit fehlenden Verknüpfungen kann dadurch steigen; sie muss aber nicht für jede Abweichung steigen. Der praktische Lernkreislauf lautet:
\begin{enumerate}
\item Eine fachlich konkrete Frage an die KI stellen.
\item Eine Antwort mit Definition, Quelle und Datenstichtag erhalten.
\item Antwort mit der eigenen Strategie oder Controlling-Analyse vergleichen.
\item Bei Abweichungen die Ursachen prüfen, einschließlich möglicher Lücken in der eigenen Analyse.
\item Datenquelle, Verknüpfung oder Definition verbessern und die Änderung dokumentieren.
\item Die nächste Antwort und die zugehörigen Testfragen erneut prüfen.
\end{enumerate}
Die Gegenprüfung ist damit zugleich Schutzmaßnahme und Diagnoseinstrument: Die Anwendung kann schon während der Verbesserung helfen und prüft gleichzeitig, ob das Datennetz die gestellte Frage trägt.

\begin{satz}[VOI-Bedingung für Datenerweiterung]
Sei $C_{int}$ die Integrationskosten einer zusätzlichen Quelle, und $\Delta W(\theta)$ der erwartete Margingewinn, der sich aus der verbesserten Antwort bei Zustand $\theta$ ergibt. Dann ist das Erweitern der Datenquelle rational genau dann, wenn
\[
\E_{\pi_1}[\Delta W(\theta)] \ge C_{int}.
\]
\end{satz}
\begin{proof}
Kostennutzen-Argument; Standard in VOI-Theorie.
\end{proof}

Damit ist klar: Eine Abweichung kann (und soll) ökonomisch als Diagnose genutzt werden. Sie ist nicht nur ein Fehler: sie zeigt, wo Investitionen in Daten den größten Margennutzen liefern.

\section{Das Data Product als gemeinsame Grundlage}
\label{sec:data-product}

Das Data Product ist die im Webinar vorgeschlagene Antwort auf die Frage, wo fachliches Wissen, relevante Daten und ihre Verbindungen so abgelegt werden, dass Menschen und KI sie wiederverwenden können. Es bündelt Daten, Bedeutung und ein Nutzungslabel. Im Beispiel könnte ein Data Product „Sales Result 360“ Fragen zu Umsatz, Marge und deren Einflussfaktoren beantworten.

\subsection{Daten und Granularität}
Im Zentrum stehen Auftragspositionen mit Umsatz, Herstellkosten, Rabatt und Menge. Die Granularität muss ausdrücklich festgelegt werden, beispielsweise „eine Zeile entspricht einer Auftragsposition“. Daran werden Kunden-, Produkt-, Zeit- und Vertriebsdimensionen sowie Fracht je Lieferung, Planwerte und Vorjahreswerte angeschlossen. Relevante Dokumente, etwa Rahmenvereinbarungen im PDF-Format, können über identifizierbare Referenzen ergänzt werden. Die Verknüpfung ist entscheidend: Ein Frachtdatensatz ohne Beziehung zur verursachenden Lieferung ist für eine Auftragsmargenfrage nur eingeschränkt nützlich.

\subsection{Bedeutung und Regeln}
Der semantische Teil beschreibt nicht nur Spaltennamen, sondern legt fest, wie Kennzahlen zu interpretieren und zu vergleichen sind.
\begin{table}[htbp]
\centering
\small
\caption{Beispielhafte fachliche Bedeutung für Sales Result 360}
\label{tab:meaning}
\begin{tabular}{p{0.25\linewidth}p{0.65\linewidth}}
\toprule
\textbf{Element} & \textbf{Festlegung} \\
\midrule
Kennzahlen & DB\,I, DB\,II, Rabattquote und Frachtquote einschließlich Berechnungsregeln. \\
Synonyme & Der unspezifische Begriff „Marge“ bezeichnet in diesem Anwendungsbereich DB\,II. \\
Zurechnungsregel & Expressfracht wird der verursachenden Lieferung und, soweit möglich, der Auftragsposition zugeordnet. \\
Einflussgrößen & Fracht, Expresslieferungen, Rabatte und Produktmix sind für die Margenanalyse zu prüfen. \\
Vergleiche & Vorjahres- und Planvergleich verwenden jeweils dokumentierte Zeiträume und Abgrenzungen. \\
\bottomrule
\end{tabular}
\end{table}
Eine dokumentierte Regel zur Zuordnung von Expressfracht würde im Fallbeispiel helfen, die 230\,T€ als Kosten der betroffenen Aufträge sichtbar zu machen. Semantik ersetzt dabei keine belastbare Verknüpfung der Quelldaten; beide Voraussetzungen müssen erfüllt sein.

\subsection{Label und Wiederverwendung}
Das Label macht Nutzungsvoraussetzungen und Grenzen sichtbar. Ein Beispiel enthält einen benannten fachlichen Kontakt, die Vollständigkeit kritischer Attribute (etwa Produktgruppe mindestens 99\,\%, täglich geprüft), Aktualisierungszusagen (etwa täglich für Vertrieb und wöchentlich für Fracht), rollenbasierten Zugriff einschließlich KI-Zugriff (beispielsweise über MCP), eine Versionsnummer wie 1.2 sowie einen Prozess zur Ankündigung von Änderungen. Ergänzend sollten Datenstand, bekannte Einschränkungen und Zugriffsbedingungen maschinenlesbar dokumentiert sein.

Ein einmal aufgebautes Data Product kann mehreren Anwendungen dienen: Controlling kann etwa mit Qlik Answers oder Copilot in Power BI nach Ursachen einer monatlichen DB-II-Veränderung fragen, das Management in Teams nach Frachtkostentrends und der Vertrieb nach Kunden mit hohem Umsatz bei geringer Marge. Ein Agent kann – beispielsweise über MCP und nur im freigegebenen Rahmen – ungewöhnliche Aufträge suchen und mögliche Ursachen erläutern. Gemeinsame Definitionen, Rechte und Protokollierung vermeiden, dass jede Anwendung dieselbe Geschäftslogik neu aufbaut oder abweichend interpretiert. Der Aufwand für verlässliche Verknüpfungen wird so einmal investiert und kann in mehreren Anwendungen Nutzen stiften.

\begin{quote}
\textbf{Faustregel:} ein Geschäftsbereich, eine Granularität, ein verantwortlicher Owner, viele Fragestellungen.
\end{quote}
Ein eng umrissenes Data Product mit wenigen zentralen Treibern ist meist ein besserer Startpunkt als der Versuch, sofort ein unternehmensweit vollständiges Modell zu bauen.

\section{Freigabe, Vorgehen und Stolperfallen}
\label{sec:release}

\subsection{Verantwortung und Freigabe}

Für einen verantwortbaren Go-live braucht es drei explizit benannte Zuständigkeiten. Die Rollen können in kleineren Unternehmen von derselben Person wahrgenommen werden; die Entscheidungen und Verantwortungsbereiche müssen trotzdem klar dokumentiert sein.

\begin{table}[htbp]
\centering
\small
\caption{Drei Verantwortlichkeiten vor dem Go-live}
\label{tab:accountabilities}
\begin{tabular}{p{0.20\linewidth}p{0.29\linewidth}p{0.22\linewidth}p{0.22\linewidth}}
\toprule
\textbf{Rolle} & \textbf{Verantwortung} & \textbf{Typische Verankerung} & \textbf{Leitfrage} \\
\midrule
Data Owner & Definiert fachliche Bedeutung, Kennzahlen, Ausnahmen und Kriterien einer richtigen Antwort; im Alltag gegebenenfalls unterstützt durch einen Data Steward. & Fachbereich, etwa Leitung Controlling. & Wer zeichnet die Definition von „Marge“ ab? \\
Data Product Owner & Liefert und verantwortet Vollständigkeit, Aktualität, Prüfungen und Berechtigungen des Data Products. & Datenteam oder IT. & Wer reagiert, wenn Frachtangaben fehlen? \\
Use-Case Owner & Genehmigt den KI-Einsatz, verantwortet Testfragen und Error Budget und legt fest, wann ein Mensch entscheidet. & Fachbereich gemeinsam mit IT. & Wer darf aus einer Antwort eine Handlung machen? \\
\bottomrule
\end{tabular}
\end{table}

Ohne Data Owner können verschiedene Werkzeuge unter „Marge“ unterschiedliche Kennzahlen verstehen. Ohne Data Product Owner bleiben ausgefallene oder verspätete Datenlieferungen womöglich unbemerkt. Ohne Use-Case Owner sind weder Testumfang noch Autonomiegrenzen eindeutig. Wenn sich für eine dieser Fragen keine verantwortliche Person benennen lässt, ist die Freigabe des Anwendungsfalls zurückzustellen.

\subsection{Vorgehen}

Ein praktikabler Aufbau beginnt mit einer geschäftlich wichtigen, aber begrenzten Frage und erweitert Netz und Nutzung erst anhand von Beobachtungen:
\begin{enumerate}
\item \textbf{Zweck und Stakes klären:} Welche Entscheidung soll unterstützt werden, welche Folgen hätte ein Fehler, und welches Error Budget ist dafür vertretbar?
\item \textbf{Einflussgrößen benennen:} Fachleute identifizieren die tatsächlichen Treiber der Antwort. Im Margenfall gehören dazu nicht nur Rabatte, sondern auch Fracht und Produktmix.
\item \textbf{Datennetz prüfen:} Für jeden Treiber wird festgestellt, ob Daten vorhanden, zugänglich, aktuell und miteinander verknüpft sind.
\item \textbf{Bedeutung festlegen:} Der Data Owner definiert Kennzahlen, Synonyme, Zurechnungsregeln und Vergleichszeiträume und zeichnet die Definition ab.
\item \textbf{Data Product aufbauen:} Daten, Granularität, Meaning und Label werden zusammengeführt; Owner, Qualität, Aktualisierung, Zugriff und Version sind dokumentiert.
\item \textbf{Referenzfragen pflegen:} Der Use-Case Owner erstellt Testfragen mit bekannten Antworten, etwa anhand einer durch Controlling geprüften Ursachenanalyse.
\item \textbf{Als Hinweis und Gegenprüfung starten:} KI-Antworten werden mit eigener Analyse und Strategie verglichen; Abweichungen erhalten eine dokumentierte Ursachenprüfung.
\item \textbf{Vertrauensstufe nur bei Nachweis erhöhen:} Von Hinweisen über Gegenprüfungen zur Entscheidungsvorbereitung; Automatisierung nur mit expliziter Freigabe, Protokollierung und Eingriffsmöglichkeit.
\item \textbf{Regelmäßig nachjustieren:} Datenqualität, Aktualisierung, Definitionen und Testfragen werden erneut geprüft. Änderungen am Data Product werden versioniert und angekündigt.
\end{enumerate}

\subsubsection{Mittelstand und Großunternehmen}
\label{sec:company-size}
Das wirtschaftliche Prinzip bleibt unabhängig von der Unternehmensgröße gleich: Die Stufe, auf der eine KI-Anwendung betrieben werden darf, richtet sich nach den Stakes der Entscheidung, nach dem Error Budget und nach den vorliegenden Nachweisen (Abschnitt~\ref{sec:error}), nicht nach der Größe der Organisation. Was sich mit der Größe ändert, sind die Ausgangslage der Datenlandschaft, der sinnvolle Zuschnitt des ersten Data Products, die Aufteilung der drei Verantwortlichkeiten, das typische Hauptrisiko und die Hebel, mit denen sich Nutzen erzielen lässt. Tabelle~\ref{tab:company-size} stellt diese fünf Aspekte gegenüber; die anschließenden Absätze erläutern jede Zeile im Zusammenhang mit dem formalen Modell und der Fallstudie.

Die Gegenüberstellung ist bewusst idealtypisch. Sie beruht auf konzeptionellen Überlegungen dieser Arbeit und ist keine empirisch geprüfte Klassifikation; eine feste Schwelle (etwa nach Mitarbeiterzahl oder Umsatz), ab der ein Unternehmen „Mittelstand“ oder „Großunternehmen“ im Sinne der Tabelle ist, wird nicht behauptet. Reale Organisationen liegen häufig dazwischen: Ein mittelständischer Zulieferer mit mehreren Standorten kennt ähnliche Silos wie ein Konzern, und ein kleiner Geschäftsbereich eines Konzerns kann wie ein Mittelständler arbeiten. Maßgeblich ist daher nicht das Etikett, sondern die Frage, welche der beiden Spalten die eigene Situation für einen konkreten Anwendungsfall besser beschreibt.

\begin{table}[H]
\centering
\small
\caption{Schwerpunktunterschiede bei der Umsetzung}
\label{tab:company-size}
\begin{tabular}{P{0.18\linewidth}P{0.36\linewidth}P{0.36\linewidth}}
\toprule
\textbf{Aspekt} & \textbf{Mittelstand} & \textbf{Großunternehmen} \\
\midrule
Ausgangslage & Überschaubare Landschaft, oft wenige Kernsysteme und knappe Spezialressourcen. & Viele Systeme, Geschäftsbereiche und Standorte mit historisch gewachsenen Silos. \\
Start & Ein Bereich, ein Data Product und wenige entscheidende Fragen. & Mehrere domänenspezifische Data Products, abgestimmt über gemeinsame Definitionen. \\
Verantwortung & Rollen können in Personalunion liegen; Zuordnung dennoch schriftlich festhalten. & Rollen trennen und durch Data Stewards im Alltag unterstützen. \\
Hauptrisiko & Datenlücken bleiben unentdeckt, wenn niemand systematisch gegenprüft. & Widersprüchliche Begriffe wie „Marge“ führen zu inkonsistenten Antworten zwischen Bereichen. \\
Hebel & Schneller Nutzen durch Abgleich mit der Controlling-Analyse; schlanke, aber explizite Governance. & Wiederverwendung, standardisierte Rechte und Protokolle sowie kontrollierte Freigaben für mehrere Anwendungen. \\
\bottomrule
\end{tabular}
\end{table}

\paragraph{Ausgangslage.}
Im Mittelstand besteht die Datenlandschaft häufig aus wenigen Kernsystemen, etwa einem ERP-System, einer Vertriebs- und einer Logistiklösung sowie ergänzenden Tabellenkalkulationen. Das hat zwei Seiten. Einerseits sind die Wege kurz: Wer im Controlling arbeitet, kennt oft informell, welche Zahl aus welchem System stammt und wie Fracht, Aufträge und Kunden zusammenhängen. Andererseits liegt dieses Wissen häufig in Köpfen statt in dokumentierten Verknüpfungen, und Spezialressourcen für Datenintegration sind knapp. In der Fallstudie zeigt sich genau dieses Muster: Die Expressfracht war im Logistiksystem gebucht, aber nicht mit den verursachenden Aufträgen verbunden, und die Produktgruppe fehlte im zugänglichen Datennetz. Beides war den Fachleuten vermutlich bekannt, für die KI aber unsichtbar. Im Sinne des formalen Modells sind die Knoten des Netzes überschaubar, doch notwendige Kanten fehlen, weil sie nur implizit existieren.

In Großunternehmen ist die Lage eine andere. Viele Systeme, Geschäftsbereiche und Standorte haben sich historisch unabhängig voneinander entwickelt; dasselbe Konzept wird in verschiedenen Silos unterschiedlich benannt, abgegrenzt und gepflegt. Das Datennetz $G$ ist groß, die Knoten unterscheiden sich stark in Qualität $q_v$ und Aktualität $t_v$, und der Nachweis, dass zwei Datenbestände verbunden sind (Abschnitt~\ref{sec:operationalization}), ist über Systemgrenzen hinweg aufwendiger zu führen. Das Problem ist weniger, dass niemand die Zusammenhänge kennt, sondern dass niemand sie vollständig kennt und dass sie sich je nach Bereich unterscheiden.

\paragraph{Start.}
Für den Mittelstand empfiehlt sich der Einstieg über einen Geschäftsbereich, ein Data Product und wenige entscheidende Fragen, im Beispiel „Sales Result 360“ mit der Frage nach den Ursachen einer sinkenden DB-II-Marge. Das entspricht der Faustregel aus Abschnitt~\ref{sec:data-product}: ein Geschäftsbereich, eine Granularität, ein verantwortlicher Owner, viele Fragestellungen. Ein eng geschnittenes Produkt lässt sich mit begrenztem Personal aufbauen, liefert früh überprüfbare Ergebnisse und erzeugt die Erfahrung, die für weitere Produkte nötig ist. Eine vollständige technische Inventur aller Systeme würde dagegen Ressourcen binden, bevor ein erster Nutzen sichtbar ist.

Großunternehmen werden typischerweise mehrere domänenspezifische Data Products aufbauen, etwa für Vertrieb, Logistik und Produktion. Damit diese Produkte zusammen Antworten tragen, müssen sie sich auf gemeinsame Definitionen stützen oder ihren jeweiligen Geltungsbereich ausdrücklich dokumentieren. Der Zuschnitt sollte sich an fachlichen Domänen und an der Granularität der Fragen orientieren und nicht allein an der eingesetzten Technologie. Der Versuch, in einem Schritt ein unternehmensweit vollständiges Modell zu bauen, verschiebt den Nutzen in die ferne Zukunft und scheitert häufig an den Abstimmungskosten.

\paragraph{Verantwortung.}
Die drei Rollen aus Tabelle~\ref{tab:accountabilities} können in kleinen Organisationen in Personalunion wahrgenommen werden. Das ist zulässig, sofern die Zuordnung schriftlich festgehalten wird. Ein Beispiel ist eine Controllerin, die zugleich Data Owner ist und die Referenzfragen als Use-Case Owner verantwortet. In diesem Fall entsteht jedoch ein Interessenkonflikt: Wer die Definition von „Marge“ festlegt und zugleich beurteilt, ob die KI-Antwort sie richtig anwendet, prüft sich teilweise selbst. Vor allem auf den Stufen~3 und~4 empfiehlt sich deshalb, die Entscheidungen zumindest gedanklich und in der Dokumentation zu trennen und, wo möglich, eine zweite Person hinzuzuziehen. Im Verfahren der Operationalisierung führt eine nicht benannte Rolle ($\rho_r=\bot$) zu einem Freigabehindernis; eine Personalunion hingegen nicht, weil die Benennung genügt.

In Großunternehmen werden die Rollen getrennt und im Alltag durch Data Stewards unterstützt. Das erhöht die Sorgfalt und die Skalierbarkeit, schafft aber ein anderes Risiko, nämlich diffuse Verantwortung: Wenn mehrere Bereiche an einem Data Product beteiligt sind, kann jeder annehmen, ein anderer reagiere auf fehlende Frachtangaben oder veraltete Daten. Klare Zuständigkeiten, Vertretungsregeln und ein Eskalationsweg sind deshalb wichtiger als zusätzliche Gremien.

\paragraph{Hauptrisiko.}
Im Mittelstand besteht das Hauptrisiko darin, dass Datenlücken unentdeckt bleiben, wenn niemand systematisch gegenprüft. Weil Prozesse kurz sind und die Beteiligten sich kennen, entsteht leicht der Eindruck, es sei bekannt, was die Daten hergeben. Die Fallstudie zeigt das Gegenteil: Die KI-Antwort erklärte nur etwa 14\,\% der Abweichung, und ohne den Abgleich mit der Controlling-Analyse wäre das nicht aufgefallen. Die systematische Gegenprüfung (Abschnitt~\ref{sec:opportunity}) ist daher im Mittelstand weniger ein Zusatz als die wichtigste Schutzmaßnahme.

In Großunternehmen steht ein anderes Risiko im Vordergrund: Widersprüchliche Begriffe führen zu inkonsistenten Antworten zwischen Bereichen. Ein Begriff wie „Marge“ kann den Deckungsbeitrag~I, den Deckungsbeitrag~II oder eine konzernweit definierte Kennzahl meinen. Im Fehlermodell entspricht das dem Bias aus uneinheitlicher Definition $B_{def}$; die Daten selbst können korrekt sein, und die Antworten weichen dennoch voneinander ab. Hinzu kommt, dass solche Inkonsistenzen mit der Zahl der beteiligten Bereiche und der Anwendungen wachsen, die dieselben Begriffe verwenden.

\paragraph{Hebel.}
Der wichtigste Hebel des Mittelstands ist der schnelle Nutzen durch den Abgleich mit der bestehenden Controlling-Analyse, verbunden mit einer schlanken, aber expliziten Governance. Schlank heißt, dass ein kurzes Label, wenige belastbare Referenzfragen und benannte Personen genügen können, solange die Anwendung auf den Stufen~1 und~2 bleibt, also als Hinweis und Gegenprüfung dient. Explizit heißt, dass auch die kleine Lösung ihre Definitionen, Verantwortlichen und Grenzen schriftlich ausweist.

Der Hebel großer Organisationen liegt in der Wiederverwendung: Der Aufwand für verlässliche Verknüpfungen, Definitionen und Rechte wird einmal investiert und kann mehreren Anwendungen dienen. Standardisierte Rechte und Protokolle gelten dann auch für KI-Agenten, etwa wenn sie über MCP auf ein Data Product zugreifen. Eine höhere Autonomiestufe wird nicht pauschal für das Unternehmen, sondern für jeden Anwendungsfall anhand seiner Stakes und seiner Testergebnisse gesondert freigegeben. Damit wird verhindert, dass ein gelungener Einsatz mit niedrigen Stakes stillschweigend auf einen Anwendungsfall mit hohen Stakes übertragen wird.

\paragraph{Übergang und Grenzen.}
Die beiden Spalten beschreiben keine getrennten Welten, sondern Pole eines Kontinuums. Ein Mittelständler, der sein erstes Data Product erfolgreich betreibt, wird früher oder später ein zweites anschließen. Spätestens wenn mehrere Anwendungen dasselbe Produkt nutzen oder erstmals Definitionskonflikte zwischen Bereichen auftreten, werden Elemente der rechten Spalte sinnvoll, etwa ein Steward für die laufende Pflege oder ein gemeinsames Verzeichnis von Synonymen. Umgekehrt kann ein Großunternehmen in einem Pilotbereich zunächst wie ein Mittelständler beginnen. Die Arbeit gibt dafür keine Schwelle vor; die Tabelle ist eine Entscheidungshilfe, die die verantwortliche Person auf die eigene Organisation anwenden muss.

\subsection{Stolperfallen}
\label{sec:pitfalls}

Auch ein sorgfältig aufgebautes Data Product schützt nicht vor typischen Fehlannahmen im Umgang mit KI-Antworten. Die Fallstudie hat bereits mehrere davon gezeigt: Eine flüssige Antwort wurde als vollständige Erklärung gelesen, die Ursache wurde zunächst im Modell statt im Datennetz vermutet, und die größten Treiber blieben unsichtbar. Dieser Abschnitt ordnet acht häufige Stolperfallen und beschreibt für jede, wie sie entsteht, welche Folge sie hat und welche Gegenmaßnahme sich aus dem Konzept der Arbeit ergibt. Die Auswahl ist eine konzeptionelle Zusammenstellung und erhebt keinen Anspruch auf Vollständigkeit oder empirische Häufigkeit.

Die Stolperfallen lassen sich drei Gruppen zuordnen. Die erste betrifft die \emph{Beurteilung von Antworten}: flüssige Antworten mit korrekten verwechseln (Nr.~1), sofort das Modell verantwortlich machen (Nr.~2) und die KI als alleinige Autorität einsetzen (Nr.~3). Die zweite betrifft die \emph{Datengrundlage und ihre Pflege}: Kennzahlen nicht definieren (Nr.~4), relevante Daten unsichtbar lassen (Nr.~5) und keine Testfragen pflegen (Nr.~6). Die dritte betrifft \emph{Tempo und Autonomie} und besteht aus zwei entgegengesetzten Fehlern: auf Perfektion warten (Nr.~7) und die Autonomie zu früh erhöhen (Nr.~8). Tabelle~\ref{tab:pitfalls} fasst alle acht Punkte in der Reihenfolge dieser Gruppen zusammen; die folgenden Absätze erläutern sie einzeln.

\begin{table}[!t]
\centering
\small
\caption{Häufige Fehlannahmen und geeignete Gegenmaßnahmen}
\label{tab:pitfalls}
\begin{tabular}{P{0.05\linewidth}P{0.21\linewidth}P{0.29\linewidth}P{0.33\linewidth}}
\toprule
\textbf{Nr.} & \textbf{Stolperfalle} & \textbf{Mögliche Folge} & \textbf{Gegenmaßnahme} \\
\midrule
\multicolumn{4}{l}{\emph{Beurteilung von Antworten}} \\
1 & Flüssige Antworten mit korrekten Antworten verwechseln & Plausible, aber unvollständige Aussagen werden ungeprüft übernommen. & Definition, Quelle und Datenstichtag einfordern. \\
2 & Sofort das Modell verantwortlich machen & Wechsel des Modells lässt fehlende Quellen und Verknüpfungen bestehen. & Abweichungen zuerst im Datennetz und in den Definitionen untersuchen. \\
3 & KI als alleinige Autorität einsetzen & Eine einzelne Antwort löst eine ungeprüfte Entscheidung aus. & Eigene Analyse gegenprüfen, insbesondere bei hohen Stakes. \\
\midrule
\multicolumn{4}{l}{\emph{Datengrundlage und Pflege}} \\
4 & Kennzahlen nicht definieren & Personen und Werkzeuge berechnen „Marge“ unterschiedlich. & Data Owner benennt und genehmigt Definitionen sowie Synonyme. \\
5 & Relevante Daten bleiben unsichtbar & Fachbereichsübergreifende Daten gehen der Analyse verloren. & Quellen und Beziehungen in das Data Product aufnehmen und Verantwortung zuweisen. \\
6 & Keine Testfragen pflegen & Qualitätsabfall bleibt unbemerkt. & Use-Case Owner hält Referenzfragen und erwartete Antworten aktuell. \\
\midrule
\multicolumn{4}{l}{\emph{Tempo und Autonomie}} \\
7 & Auf Perfektion warten & Nützliche, risikoarme Anwendungen werden unbegrenzt vertagt. & Mit klar begrenzten Hinweisen und Gegenprüfungen beginnen. \\
8 & Autonomie zu früh erhöhen & Unvollständige Informationen lösen Maßnahmen ohne menschliche Prüfung aus. & Freigabe, Protokollierung und menschliche Eingriffspunkte vorschreiben. \\
\bottomrule
\end{tabular}
\end{table}

\subsubsection{Beurteilung von Antworten}

\paragraph{Nr.~1: Flüssige Antworten mit korrekten verwechseln.}
Sprachmodelle formulieren auch dann geschlossen und selbstsicher, wenn ihre Datengrundlage lückenhaft ist. Sprachliche Qualität ist deshalb kein Indikator für inhaltliche Vollständigkeit. In der Fallstudie klang die Empfehlung zu Rabattgenehmigungen und Neuverhandlungen plausibel und konkret, enthielt aber weder eine Quellenangabe noch einen Datenstichtag und erklärte nur rund 14\,\% der Abweichung. Wer die Antwort nur nach ihrem Eindruck beurteilt, hätte den Mangel nicht bemerkt. Die Gegenmaßnahme besteht darin, Definition, Quelle und Datenstichtag systematisch einzufordern; das entspricht der Provenance aus Abschnitt~\ref{sec:formal}. Messbar wird sie über die Provenance-Abdeckung in Tabelle~\ref{tab:metrics}, und im Prüfverfahren hängt der Nachweis $e_{\mathrm{prov}}$ daran, dass jede Antwort eine vollständige Provenance trägt (Definition~\ref{def:evidence}).

\paragraph{Nr.~2: Sofort das Modell verantwortlich machen.}
Wenn eine Antwort enttäuscht, liegt die Vermutung nahe, das Modell sei zu schwach. Diese Annahme führt in die Irre, wenn die Ursache im Netz liegt. Fehlen Fracht-Auftrag-Verknüpfung und Produktgruppe, findet auch ein leistungsfähigeres Modell dieselben Lücken vor und kann die fehlenden Treiber nicht erfinden; ein Modellwechsel verbraucht dann Zeit und Budget, ohne den Fehler zu verändern. Der Fehleranteil aus fehlenden Einflussgrößen $B_{miss}$ und aus uneinheitlichen Definitionen $B_{def}$ (Abschnitt~\ref{sec:error}) ist vom Modell weitgehend unabhängig. Die Gegenmaßnahme ist eine feste Reihenfolge der Ursachenprüfung, wie sie Tabelle~\ref{tab:deviation-causes} beschreibt: zuerst Quellen, Attribute, Definitionen und Aktualität, dann die eigene Analyse, erst danach das Modell. Das schließt nicht aus, dass ein Modell tatsächlich schwach sein kann; es soll nur verhindern, dass diese Erklärung die einfachste Antwort ist und deshalb zuerst gewählt wird.

\paragraph{Nr.~3: KI als alleinige Autorität einsetzen.}
Wird eine einzelne Antwort ohne Gegenprüfung zur Grundlage einer Entscheidung, hängt das Risiko vollständig von der Qualität dieser einen Antwort ab. Im Beispiel hätte die Rabattverhandlung mit dem Kunden Nord, der 2{,}0\,Mio.\,€ Umsatz erwirtschaftete, eine Geschäftsbeziehung belastet, während die größte Ursache, die Expressfracht mit 230\,T€, unberührt geblieben wäre. Je höher die Stakes, desto weniger darf eine Antwort allein tragen (Tabelle~\ref{tab:stakes}). Die Gegenmaßnahme besteht darin, die Antwort mit der eigenen Analyse, der Strategie oder einer fachlich geprüften Referenz abzugleichen und einen Menschen entscheiden zu lassen, solange keine Automatisierung auf Stufe~4 ausdrücklich freigegeben ist. Auf Stufe~3 ist die Forderung, dass ein Mensch entscheidet, deshalb ein eigener Nachweis im Prüfverfahren ($e_{\mathrm{hum}}$).

\subsubsection{Datengrundlage und Pflege}

\paragraph{Nr.~4: Kennzahlen nicht definieren.}
Begriffe wie „Marge“ werden im Alltag selbstverständlich verwendet, aber nicht immer gleich verstanden. Meint der Controller den Deckungsbeitrag~II, der Vertrieb den Deckungsbeitrag~I und die KI eine aus dem Kontext geratene Größe, entsteht ein Unterschied, der keiner Datenlücke entspricht und trotzdem zu abweichenden Antworten führt. Diese Stolperfalle erzeugt den Definitionsbias $B_{def}$ und ist besonders tückisch, weil jede beteiligte Zahl für sich korrekt sein kann. Die Gegenmaßnahme ist, dass der Data Owner Kennzahlen, Synonyme und Geltungsbereich festlegt und abzeichnet (Tabelle~\ref{tab:meaning}). Im Prüfverfahren ist die Definition jedes Kennzahlbegriffs, auch über ein Synonym, Teil des Nachweises $e_{\mathrm{kpi}}$.

\paragraph{Nr.~5: Relevante Daten bleiben unsichtbar.}
Daten, die in einem anderen Fachbereich oder System liegen, gehen der Analyse verloren, wenn sie nicht ins Data Product aufgenommen und mit den übrigen Daten verknüpft werden. Das war in der Fallstudie die Hauptursache: Expressfracht und Produktmix machten zusammen rund 86\,\% der Abweichung aus und waren für die KI nicht sichtbar. Besonders schwer zu erkennen ist das, weil die vorhandenen Daten selbst unauffällig sind; das Problem ist, was fehlt. Die Gegenmaßnahme besteht darin, die wesentlichen Einflussgrößen zuvor mit Fachleuten zu benennen (Schritt~2 des Vorgehens), Quellen und Beziehungen ins Data Product aufzunehmen und die Verantwortung dafür dem Data Product Owner zuzuweisen. Das Prüfverfahren fordert dazu, dass alle geforderten Verknüpfungen nachweisbar sind, sonst bleibt $e_{\mathrm{kpi}}=0$.

\paragraph{Nr.~6: Keine Testfragen pflegen.}
Ein Data Product wird nach dem Go-live verändert: Quellsysteme werden angepasst, Definitionen verfeinert, Modelle ausgetauscht und Berechtigungen geändert. Ohne gepflegte Referenzfragen mit erwarteten Antworten bleibt unbemerkt, wenn die Qualität nachlässt. Die Gefahr ist schleichend, weil einzelne Antworten weiterhin plausibel wirken. Als Gegenmaßnahme hält der Use-Case Owner die Referenzfragen aktuell, gewichtet sie nach Stakes und führt sie nach jeder relevanten Änderung erneut aus (Abschnitt~\ref{sec:metrics}). Zwei Grenzen sind zu beachten: Ein hoher Wert auf leicht beantwortbaren Fragen sagt wenig über seltene, aber folgenreiche Ausnahmen aus, und ein Vergleich über die Zeit ist nur aussagekräftig, wenn der Testsatz unverändert blieb (Bemerkung zu den Grenzen des Vergleichs).

\subsubsection{Tempo und Autonomie}

\paragraph{Nr.~7: Auf Perfektion warten.}
Die Forderung nach vollständig fehlerfreien Daten ist praktisch unerreichbar und, wie die ökonomische Gut-genug-Bedingung zeigt, auch nicht nötig. Wer den Einsatz bis dahin verschiebt, verzichtet auf Hinweise und Gegenprüfungen, die schon bei geringer Abdeckung nützlich sind; im Margenbeispiel war der Rabatthinweis wahr und prüfenswert, obwohl er nur einen kleinen Teil erklärte. Zudem entgeht dem Unternehmen der Lerneffekt: Abweichungen zeigen, wo Investitionen in Daten den größten Nutzen bringen (Abschnitt~\ref{sec:opportunity}). Die Gegenmaßnahme ist, mit klar begrenzten Anwendungen auf den Stufen~1 und~2 zu beginnen, ihre Grenzen offen auszuweisen und die Stufe erst bei belegtem Reifegrad zu erhöhen.

\paragraph{Nr.~8: Autonomie zu früh erhöhen.}
Das Gegenstück zum Warten ist die voreilige Automatisierung. Läuft eine Anwendung über einige Zeit unauffällig, entsteht der Wunsch, sie ohne menschliche Prüfung Maßnahmen auslösen zu lassen. Fehlen dabei Informationen, wie sie in der Fallstudie fehlten, werden Maßnahmen auf unvollständiger Grundlage ausgeführt, und der Schaden tritt ein, bevor ihn jemand bemerkt. Eine höhere Stufe folgt nicht daraus, dass die vorherige eine Zeit lang funktioniert hat, sondern setzt einen belegbaren Reifegrad voraus (Abschnitt~\ref{sec:error}). Die Gegenmaßnahme ist, für Stufe~4 ausdrückliche Freigabe, Protokollierung und menschliche Eingriffspunkte vorzuschreiben; im Prüfverfahren sind das die Nachweise $e_{\mathrm{app}}$, $e_{\mathrm{log}}$ und $e_{\mathrm{int}}$.

\subsubsection{Wechselwirkungen und Frühwarnzeichen}
Die acht Stolperfallen treten selten einzeln auf. Wer flüssige Antworten für korrekt hält (Nr.~1) und die KI zugleich als Autorität ansieht (Nr.~3), ist besonders anfällig dafür, die Autonomie zu früh zu erhöhen (Nr.~8). Nicht definierte Kennzahlen (Nr.~4) und unsichtbare Daten (Nr.~5) erzeugen Abweichungen, die ohne Testfragen (Nr.~6) niemand systematisch auswertet, und wer dann vorschnell das Modell verantwortlich macht (Nr.~2), schließt die Lücke nicht. Die Nummern~7 und~8 bilden ein Spannungsfeld: Wer zu lange wartet, verschenkt Nutzen; wer zu früh automatisiert, trägt unnötige Risiken. Das Stufenmodell der Arbeit (Tabelle~\ref{tab:reliability-levels}) löst diese Spannung, indem es mit niedrigem Risiko beginnt und jeden weiteren Schritt an Nachweise bindet.

Als Frühwarnzeichen eignen sich Beobachtungen, die sich ohne großen Aufwand erheben lassen: Antworten ohne Quellenangabe oder Stichtag, ein Modellwechsel als Reaktion auf eine Abweichung, Begriffe ohne benannte Verantwortliche, Testfragen mit veraltetem Datum, Anwendungen, die ohne dokumentierte Freigabe Handlungen auslösen, sowie Projekte, deren Start an dem Wunsch nach vollständigen Daten hängt. Wer eines oder mehrere dieser Zeichen beobachtet, sollte die Checkliste in Abschnitt~\ref{sec:checklist} durchgehen und die Freigabe des Anwendungsfalls zurückstellen, bis die Ursache geklärt ist.

\FloatBarrier

\section{Fortschritt messen}
\label{sec:metrics}

Die Frage, ob ein Datennetz besser wird, braucht Messgrößen, die sowohl Datenzustand als auch tatsächlichen Anwendungserfolg erfassen. Einzelne Qualitätswerte ersetzen dabei keine Überprüfung der wirtschaftlich relevanten Treiber.

\begin{table}[htbp]
\centering
\small
\caption{Kennzahlen zur Beobachtung von Reife und Nutzen}
\label{tab:metrics}
\begin{tabular}{p{0.30\linewidth}p{0.60\linewidth}}
\toprule
\textbf{Kennzahl} & \textbf{Beantwortete Frage} \\
\midrule
Erklärter Anteil & Welchen Anteil eines bekannten Effekts erklärt die KI auf Referenzfällen? Im Fallbeispiel waren es etwa 14\,\%. \\
Bestehensquote der Testfragen & Welcher Anteil der gepflegten Fragen wird innerhalb der festgelegten Toleranz korrekt beantwortet? \\
Vollständigkeit kritischer Attribute & Wie vollständig sind entscheidende Felder, beispielsweise Produktgruppen mit mindestens 99\,\% Belegung? \\
Aktualität & Werden Daten innerhalb der zugesagten Intervalle aktualisiert, etwa täglich oder wöchentlich? \\
Provenance-Abdeckung & Bei welchem Anteil der Antworten sind Definition, Quelle und Datenstichtag verfügbar? \\
Bearbeitete Abweichungen & Wie viele Abweichungen führten zu einer dokumentierten Korrektur oder zu einer begründeten Bestätigung der bisherigen Daten? \\
Wiederverwendung & Wie viele Anwendungen nutzen dieselbe freigegebene Daten- und Definitionsgrundlage? \\
\bottomrule
\end{tabular}
\end{table}

Der erklärte Anteil ist eine nützliche Kennzahl, aber kein universeller Qualitätsbeweis: Ein hoher Wert in einem kleinen oder wenig repräsentativen Testsatz garantiert keine Sicherheit bei einer anderen Entscheidung. Deshalb müssen Testfälle nach Stakes gewichtet, regelmäßig aktualisiert und um typische Abweichungen ergänzt werden. Ein Netz ist für einen Anwendungsfall dann gut genug, wenn die relevanten Kennzahlen die für diesen Fall festgelegten Anforderungen erfüllen.

\subsection{Testfragen und erklärte Abdeckung}
Sei $T=\{(q_j,y_j^*)\}_{j=1}^n$ eine Menge repräsentativer Testfragen mit fachlich geprüften Referenzwerten. Für eine Fehlertoleranz $\delta_j$ und eine ökonomische Gewichtung $w_j\geq 0$ lässt sich die gewichtete Bestehensquote definieren als
\[
\mathrm{Bestehensquote}
=\frac{\sum_{j=1}^n w_j\,\mathbf{1}\!\left[|\hat y_j-y_j^*|\leq\delta_j\right]}
{\sum_{j=1}^n w_j}.
\]
Dabei bezeichnet $\mathbf{1}[\cdot]$ die Indikatorfunktion. Für Ursachenanalysen sollte zusätzlich der erklärte Anteil am bekannten Gesamteffekt ausgewiesen werden. Beide Kennzahlen beantworten verschiedene Fragen: Besteht die Antwort einen Toleranztest, und deckt sie die wesentlichen Ursachen ab? Die Referenzen und Toleranzen verantwortet der Use-Case Owner gemeinsam mit den fachlich zuständigen Personen.

\section{Operationalisierung}
\label{sec:operationalization}

Die bisherigen Abschnitte beschreiben die Kriterien für ein ausreichendes Datennetz in Worten, Tabellen und Formeln. Dieser Abschnitt macht daraus ein deterministisches Prüfverfahren: Data Product und Anwendungsfall werden maschinenlesbar spezifiziert, messbare Nachweise werden aus den Daten \emph{abgeleitet} statt erklärt, und das Ergebnis ist die Menge der Freigabehindernisse. Ein späterer Bericht lässt sich mit einem früheren vergleichen. Das Verfahren ist im Verzeichnis \texttt{src/} umgesetzt (Module \texttt{spec}, \texttt{report}, \texttt{compare}, \texttt{cli}); formuliert werden hier nur Eigenschaften, die sich beweisen lassen. Alle Größen sind exakt rational ($\QQ$); es gibt weder Zufall noch Gleitkomma noch einen Zugriff auf die Systemuhr (der Stichtag ist Eingabe).

\subsection{Spezifikation und Eingaben}

\begin{definition}[Spezifikation]\label{def:spec}
Eine \emph{Spezifikation} ist ein Tupel
\[
\Sigma=(P,R,K,\rho,s,\beta,\theta,T,o)
\]
mit
\begin{itemize}
\item einem Data Product $P=(V_P,E_P,\mathcal M,\ell)$ gemäß Abschnitt~\ref{sec:formal},
\item geforderten Verknüpfungen $R\subseteq V\times V$ und einer endlichen Menge $K$ von Kennzahlbegriffen,
\item $\rho=(\rho_d,\rho_p,\rho_u)$: für Data Owner, Data Product Owner und Use-Case Owner je eine benannte Person oder $\bot$ (nicht benannt),
\item einer angestrebten Stufe $s\in\{1,2,3,4\}$ (Tabelle~\ref{tab:reliability-levels}),
\item einem optionalen Error Budget $\beta=(W,\varepsilon)$ mit $W\ge 0$, $\varepsilon>0$ oder $\beta=\bot$ sowie einer optionalen Mindest-Bestehensquote $\theta\in[0,1]\cup\{\bot\}$,
\item Referenzfragen $T=\bigl((q_j,y_j^*,\delta_j,w_j)\bigr)_{j=1}^n$ mit paarweise verschiedenen $q_j$, $\delta_j,w_j\ge 0$ und $\sum_j w_j>0$ falls $n\ge 1$ (Abschnitt~\ref{sec:metrics}),
\item erklärten organisatorischen Nachweisen $o=(o_{\mathrm{hum}},o_{\mathrm{app}},o_{\mathrm{log}},o_{\mathrm{int}})\in\{0,1\}^4$: ein Mensch entscheidet, ausdrückliche Freigabe, Protokollierung, Eingriffspfad.
\end{itemize}
\end{definition}

Zu einem Prüflauf gehören außerdem die Antworten $a=\bigl((\hat y_j,\Pi_j)\bigr)_{j=1}^n$ mit $\hat y_j\in\QQ\cup\{\bot\}$ ($\bot$: keine Antwort) und Provenance $\Pi_j$ (oder $\bot$), ein Stichtag $\tau$ und ein erwarteter Verlust $\bar L\in\QQ_{\ge 0}\cup\{\bot\}$, also eine Schätzung von $\E[L(\hat y,y^*)]$ aus Abschnitt~\ref{sec:error}. Der Wert $\bar L$ wird nicht vom Verfahren geschätzt, sondern geliefert.

\begin{bemerkung}[Messen statt Erklären]
Kennzahlen, Provenance, Qualität und Error Budget dürfen in $\Sigma$ nicht als erfüllt erklärt werden; eine Spezifikation, die sie enthält, wird abgelehnt. Erklärbar sind nur die vier Komponenten von $o$, weil sie organisatorische Tatsachen sind, die sich nicht aus Daten ablesen lassen. Die Richtigkeit dieser Erklärungen bleibt Sache der Organisation.
\end{bemerkung}

\subsection{Abgeleitete Nachweise}

Für ein Element $(u,v)\in R$ heißt $u$ mit $v$ \emph{nachweisbar verbunden}, wenn $u,v\in V_P$ gilt und im ungerichteten Graphen $\bigl(V_P,\{\{x,y\}: (x\to y,\varphi)\in E_P\}\bigr)$ ein Weg von $u$ nach $v$ existiert. Eine Provenance $\Pi_j$ heißt \emph{vollständig}, wenn sie mindestens eine Quelle, einen Datenstichtag und mindestens eine angewendete Regel enthält (Definition der Provenance in Abschnitt~\ref{sec:formal}); fehlt $\Pi_j$ ganz, gilt sie als unvollständig. Ein Knoten $v\in V_P$ mit letzter Aktualisierung $t_v$ und Aktualisierungsfrequenz $f$ (in Tagen, Teil des Labels) heißt \emph{aktuell zum Stichtag $\tau$}, wenn $t_v\neq\bot$ und $0\le\tau-t_v\le f$ gilt; ein unbekanntes oder in der Zukunft liegendes $t_v$ ist demnach nicht aktuell. Mit
\[
c(a)=\frac{\bigl|\{\,j:\Pi_j\text{ vollständig}\,\}\bigr|}{n},\qquad
p(a)=\frac{\sum_{j=1}^n w_j\,\mathbf 1\!\left[\hat y_j\neq\bot\ \wedge\ |\hat y_j-y_j^*|\le\delta_j\right]}{\sum_{j=1}^n w_j}
\]
(für $n\ge 1$) bezeichnen $c$ die Provenance-Abdeckung und $p$ die gewichtete Bestehensquote aus Abschnitt~\ref{sec:metrics}; eine fehlende Antwort besteht den Test nie.

\begin{definition}[Evidenzvektor]\label{def:evidence}
Der \emph{Evidenzvektor} $e=(e_f)_{f\in\mathcal F}\in\{0,1\}^{\mathcal F}$ über der Indexmenge $\mathcal F=\{\mathrm{kpi},\mathrm{prov},\mathrm{qual},\mathrm{err},\mathrm{hum},\mathrm{app},\mathrm{log},\mathrm{int}\}$ ist festgelegt durch
\begin{align*}
e_{\mathrm{kpi}}=1&\iff K\neq\emptyset\ \wedge\ \text{alle }(u,v)\in R\text{ sind nachweisbar verbunden}\\
&\qquad\wedge\ \text{jedes }k\in K\text{ ist in }\mathcal M\text{ definiert (auch über ein Synonym)},\\
e_{\mathrm{prov}}=1&\iff n\ge 1\ \wedge\ c(a)=1,\\
e_{\mathrm{qual}}=1&\iff V_P\neq\emptyset\ \wedge\ q_v\neq\bot\text{ für alle }v\in V_P,\\
e_{\mathrm{err}}=1&\iff \beta\neq\bot\ \wedge\ \bar L\neq\bot\ \wedge\ \bar L\le\varepsilon W,
\end{align*}
sowie $e_{\mathrm{hum}}=o_{\mathrm{hum}}$, $e_{\mathrm{app}}=o_{\mathrm{app}}$, $e_{\mathrm{log}}=o_{\mathrm{log}}$ und $e_{\mathrm{int}}=o_{\mathrm{int}}$.
\end{definition}

Die Definition setzt den Grundsatz der Arbeit um, dass ein fehlender Nachweis selbst ein Ergebnis ist: Jede Komponente ist nur dann $1$, wenn der Nachweis tatsächlich vorliegt. Die Bedingung für $e_{\mathrm{err}}$ ist die ökonomische Gut-genug-Bedingung aus Abschnitt~\ref{sec:error}.

\subsection{Freigabehindernisse}

Die Anforderungen wachsen mit der Stufe. Sei $F_1=\emptyset$, $F_2=\{\mathrm{kpi}\}$, $F_3=\{\mathrm{prov},\mathrm{qual},\mathrm{err}\}$, $F_4=\{\mathrm{app},\mathrm{log},\mathrm{int}\}$ und $\mathrm{Req}(s)=\bigcup_{j\le s}F_j$.

\begin{definition}[Freigabehindernisse]\label{def:blockers}
Zu Spezifikation, Antworten, Stichtag und erwartetem Verlust ist die Menge $B(s)=B_1\cup B_2(s)\cup B_3(s)\cup B_4(s)\cup B_5$ der \emph{Freigabehindernisse} die (nach Art getaggte, also disjunkte) Vereinigung von
\begin{align*}
B_1&=\{\,(\mathrm{Rolle},r): \rho_r=\bot,\ r\in\{d,p,u\}\,\},\\
B_2(s)&=\{\,f\in\mathrm{Req}(s): e_f=0\,\},\\
B_3(s)&=\{\mathrm{hum}\}\ \text{falls } s=3\text{ und }e_{\mathrm{hum}}=0,\ \text{sonst }\emptyset,\\
B_4(s)&=\{\,(\mathrm{Aktualität},v): v\in V_P\text{ nicht aktuell zum Stichtag }\tau\,\}\ \text{falls }s\ge 3,\ \text{sonst }\emptyset,\\
B_5&=\{\mathrm{pass}\}\ \text{falls }\theta\neq\bot\text{ und }(n=0\text{ oder }p(a)<\theta),\ \text{sonst }\emptyset.
\end{align*}
\end{definition}

$B_3$ entspricht der Tabelle~\ref{tab:reliability-levels}: Auf Stufe~3 entscheidet ein Mensch; auf Stufe~4 entfällt diese Forderung zugunsten von Freigabe, Protokollierung und Eingriffspfad. Die leere Menge $B=\emptyset$ ist keine Freigabe, sondern nur die Feststellung, dass keine der hier formalisierten Bedingungen verletzt ist. Sie ist \emph{notwendig, nicht hinreichend}; die Freigabe verantwortet der Use-Case Owner.

\begin{satz}[Antitonie in den Nachweisen]\label{thm:antitone}
Seien $e\le e'$ komponentenweise und alle übrigen Eingaben gleich. Dann gilt $B(s;e')\subseteq B(s;e)$ für jede Stufe $s$.
\end{satz}
\begin{proof}
$B_1$, $B_4$ und $B_5$ hängen von $e$ nicht ab und stimmen für $e$ und $e'$ überein. Sei $f\in B_2(s;e')$, also $f\in\mathrm{Req}(s)$ und $e'_f=0$. Aus $e_f\le e'_f=0$ folgt $e_f=0$, also $f\in B_2(s;e)$. Sei $\mathrm{hum}\in B_3(s;e')$, also $s=3$ und $e'_{\mathrm{hum}}=0$; dann ist $e_{\mathrm{hum}}\le 0$, also $\mathrm{hum}\in B_3(s;e)$. Damit ist jede Komponente von $B(s;e')$ in der entsprechenden Komponente von $B(s;e)$ enthalten.
\end{proof}

\begin{satz}[Stufenmonotonie bis auf die Aufsicht]\label{thm:levels}
Für $s\le s'$ und gleiche Eingaben gilt $B(s)\subseteq B(s')\cup\{\mathrm{hum}\}$.
\end{satz}
\begin{proof}
$B_1$ und $B_5$ hängen nicht von der Stufe ab. Aus $s\le s'$ folgt $\mathrm{Req}(s)\subseteq\mathrm{Req}(s')$ und damit $B_2(s)\subseteq B_2(s')$. Ist $s\ge 3$, so ist auch $s'\ge 3$; also gilt $B_4(s)\subseteq B_4(s')$ (für $s<3$ ist $B_4(s)=\emptyset$). Das einzige Element, das in $B(s)$, aber nicht in $B(s')$ liegen kann, ist $\mathrm{hum}\in B_3(s)$ (nämlich für $s=3$, $s'=4$); es steht rechts in der Vereinigung.
\end{proof}

\begin{korollar}\label{cor:levels}
Ist $B(s')=\emptyset$, so gilt $B(s)\subseteq\{\mathrm{hum}\}$ für alle $s\le s'$. Eine Stufe kann also frei von Hindernissen sein, während eine niedrigere nur an der Forderung scheitert, dass ein Mensch entscheidet; das ist kein Widerspruch, sondern Folge davon, dass Stufe~4 diese Forderung durch Freigabe, Protokoll und Eingriffspfad ersetzt.
\end{korollar}
\begin{proof}
Satz~\ref{thm:levels} liefert $B(s)\subseteq B(s')\cup\{\mathrm{hum}\}=\{\mathrm{hum}\}$.
\end{proof}

\begin{satz}[Unbekanntes zählt nie als Erfolg]\label{thm:conservative}
Die Eingaben $\Phi'$ gehen aus $\Phi$ durch genau eine der folgenden Ersetzungen hervor:
(a)~$q_v\mapsto\bot$ für ein $v\in V_P$; (b)~$t_v\mapsto\bot$ für ein $v\in V_P$; (c)~$\bar L\mapsto\bot$; (d)~in einer Provenance $\Pi_j$ wird die Quelle, der Datenstichtag oder die Regel entfernt oder $\Pi_j\mapsto\bot$ gesetzt; (e)~$\hat y_j\mapsto\bot$ für ein $j$. Dann gilt $B(\Phi)\subseteq B(\Phi')$.
\end{satz}
\begin{proof}
$B$ hängt von den Eingaben nur über drei Größen ab: den Evidenzvektor $e$ (über $B_2$ und $B_3$), die Menge $\mathrm{Fr}\subseteq V_P$ der aktuellen Knoten (über $B_4=\{(\mathrm{Aktualität},v):v\in V_P\setminus\mathrm{Fr}\}$) und die Bestehensquote $p$ (über $B_5$). $B$ ist antiton in $e$ (Satz~\ref{thm:antitone}), antiton in $\mathrm{Fr}$ (eine kleinere Menge $\mathrm{Fr}$ hat ein größeres Komplement) und antiton in $p$: Die Bedingung $p(a)<\theta$ ist für kleineres $p$ leichter erfüllt, und $n=0$ hängt von den Ersetzungen nicht ab. Es genügt daher zu zeigen, dass jede Ersetzung $e$, $\mathrm{Fr}$ und $p$ höchstens verkleinert, ohne sie zu vergrößern:
\begin{itemize}
\item (a): Die Bedingung „$q_v\neq\bot$ für alle $v\in V_P$“ kann nur von wahr zu falsch wechseln, also $e'_{\mathrm{qual}}\le e_{\mathrm{qual}}$; $e_{\mathrm{kpi}},e_{\mathrm{prov}},e_{\mathrm{err}}$, $\mathrm{Fr}$ und $p$ ändern sich nicht.
\item (b): Aktuell zu sein setzt $t_v\neq\bot$ voraus; also $\mathrm{Fr}'\subseteq\mathrm{Fr}$, alles andere bleibt gleich.
\item (c): $e_{\mathrm{err}}$ verlangt $\bar L\neq\bot$; also $e'_{\mathrm{err}}\le e_{\mathrm{err}}$.
\item (d): Die Vollständigkeit von $\Pi_j$ kann nur von wahr zu falsch wechseln; damit gilt $c(a')\le c(a)$ und folglich $e'_{\mathrm{prov}}\le e_{\mathrm{prov}}$ (denn $e_{\mathrm{prov}}=1$ verlangt $c=1$).
\item (e): Der Summand $w_j\mathbf 1[\hat y_j\neq\bot\wedge\cdots]$ wird zu $0$, während der Nenner $\sum_j w_j$ unverändert bleibt; wegen $w_j\ge 0$ gilt $p(a')\le p(a)$.
\end{itemize}
$B_1$ und die Komponente $e_{\mathrm{hum}}$ sind von allen fünf Ersetzungen unberührt.
\end{proof}

\begin{satz}[Reproduzierbarkeit]\label{thm:repro}
Bericht und Fingerprint sind Funktionen der Eingaben: Zwei Prüfläufe mit denselben Eingaben (Spezifikation, Antworten, Stichtag, erwarteter Verlust, Versionen von Python und Abhängigkeiten) liefern denselben Bericht und denselben Fingerprint.
\end{satz}
\begin{proof}
Jede Größe der Definitionen~\ref{def:evidence} und~\ref{def:blockers} ist eine Funktion der Eingaben; die Rechnung ist exakt rational und verwendet weder Zufall noch Systemzeit. Der Fingerprint ist der SHA-256-Wert einer kanonischen JSON-Serialisierung (sortierte Schlüssel, Brüche als Text) von Name, Spezifikation, Antworten, Stichtag, erwartetem Verlust und Versionen; Serialisierung und Hashfunktion sind deterministisch. Das Ergebnis des Laufs geht nicht in den Fingerprint ein, ist aber nach dem Gesagten selbst durch die Eingaben festgelegt.
\end{proof}

Die Umkehrung gilt nur eingeschränkt: Gleiche Fingerprints bedeuten gleiche kanonisch serialisierte Eingaben nur unter der Annahme, dass keine SHA-256-Kollision vorliegt.

\subsection{Vergleich über die Zeit}

Für Berichte $r$ desselben Data Products besteht der Messvektor aus sieben Komponenten: der Bestehensquote $p$, der Provenance-Abdeckung $c$ und der Zahl $u$ unbeantworteter Fragen (alle $\bot$, wenn $n=0$), sowie den Anzahlen offener Verknüpfungen, nicht aktueller Knoten (einschließlich unbekannter Aktualität), Knoten ohne bekannte Qualität und Freigabehindernisse $|B|$. Auf jeder Komponente ist eine totale Ordnung $\le_i$ „schlechter oder gleich“ festgelegt: Für $p$ und $c$ gilt $\bot\le_i x$ für alle $x$, sonst die übliche Ordnung; für alle anderen Komponenten gilt ebenfalls $\bot\le_i x$ für alle $x$, und für Zahlen $x,y\in\NN$ ist $x\le_i y$ genau dann, wenn $x\ge y$. Fehlender Nachweis ist also stets am schlechtesten.

\begin{definition}[Regressionsfreiheit]\label{def:regression}
Der Übergang vom früheren Bericht $r$ zum späteren Bericht $r'$ ist \emph{regressionsfrei}, geschrieben $r\sqsubseteq r'$, wenn \emph{erstens} $m_i(r)\le_i m_i(r')$ für jede Komponente $i$ gilt und \emph{zweitens} $B(r')\subseteq B(r)$ ist, also kein Freigabehindernis neu hinzukommt. Eine \emph{Regression} liegt vor, wenn $r\sqsubseteq r'$ nicht gilt.
\end{definition}

Dass zusätzlich $B(r')\subseteq B(r)$ verlangt wird, ist nötig: Ersetzt ein neues Hindernis ein behobenes, bleibt $|B|$ gleich, die Lage hat sich aber geändert.

\begin{satz}[Quasiordnung]\label{thm:preorder}
Die Relation $\sqsubseteq$ ist reflexiv und transitiv. Für jede Komponente gilt: $r\to r'$ ist eine Verbesserung genau dann, wenn $r'\to r$ eine Verschlechterung ist. Gilt $r\sqsubseteq r'$ und $r'\sqsubseteq r$, so stimmen alle Komponenten des Messvektors und die Mengen $B(r)$, $B(r')$ überein.
\end{satz}
\begin{proof}
Reflexivität: $m_i(r)\le_i m_i(r)$, weil $\le_i$ eine Ordnung ist, und $B(r)\subseteq B(r)$. Transitivität: Aus $r\sqsubseteq r'$ und $r'\sqsubseteq r''$ folgt $m_i(r)\le_i m_i(r')\le_i m_i(r'')$, also $m_i(r)\le_i m_i(r'')$, sowie $B(r'')\subseteq B(r')\subseteq B(r)$. Spiegelung: Weil $\le_i$ total ist, bedeutet „Verbesserung“ $m_i(r)<_i m_i(r')$ und „Verschlechterung“ $m_i(r')<_i m_i(r)$; vertauscht man $r$ und $r'$, vertauschen sich die beiden Aussagen. Antisymmetrie: Aus $m_i(r)\le_i m_i(r')\le_i m_i(r)$ folgt $m_i(r)=m_i(r')$, aus $B(r')\subseteq B(r)\subseteq B(r')$ folgt $B(r)=B(r')$.
\end{proof}

\begin{bemerkung}[Grenzen des Vergleichs]
Der Vergleich ist exakt und kennt keine Toleranz; Messrauschen (etwa bei kleinen Testmengen) modelliert er nicht. Außerdem ist $p$ nur vergleichbar, wenn die Referenzfragen gleich geblieben sind; sonst vermischt sich eine Änderung des Datennetzes mit einer Änderung des Testsatzes. Das Verfahren weist deshalb aus, ob sich die Eingaben (der Fingerprint) geändert haben, kann die beiden Ursachen aber nicht trennen.
\end{bemerkung}

\subsection{Beispiel und Grenzen}

Für das Data Product „Sales Result 360“ der Fallstudie (Datei \texttt{examples/\allowbreak sales\_\allowbreak result\_\allowbreak 360.yaml}, illustrativ) liefert das Verfahren auf Stufe~3 mit $\bar L=380$ und $\varepsilon W=22$ sieben Freigabehindernisse: vier fehlende Nachweise (Kennzahlen nicht verknüpft, weil die Verbindung zwischen Auftrag und Fracht fehlt; Provenance unvollständig, weil die Antwort zum Produktmix fehlt; Qualität des Knotens ohne Angabe nicht messbar; Error Budget überschritten), zwei nicht aktuelle Knoten und eine Bestehensquote $p=1/6<\theta=3/4$. Das Ergebnis deckt sich mit der Aussage der Fallstudie, dass das Netz für diese Frage nicht ausreicht.

Das Verfahren prüft nur, was in Definition~\ref{def:blockers} formalisiert ist. Es kontrolliert nicht, ob erklärte organisatorische Nachweise zutreffen, ob eine Quelle inhaltlich stimmt (Vollständigkeit der Provenance ist eine formale Eigenschaft), ob die gemeldeten Aktualisierungszeitpunkte korrekt sind und wie $\bar L$ geschätzt wurde. Die Stufen bleiben ein konzeptioneller Rahmen und keine Norm. Wie das Verfahren als Dienst bereitgestellt und in Abläufe eingebunden werden kann, behandelt Abschnitt~\ref{sec:api}.

\section{Bereitstellung als JSON-API}
\label{sec:api}

Die bisherigen Abschnitte haben das Prüfverfahren als Bibliothek und als Kommandozeile beschrieben (Abschnitt~\ref{sec:operationalization}). Beides setzt voraus, dass der Anwender Python-Code ausführt oder eine Shell bedient. Datennetze und KI-Anwendungen leben jedoch in Umgebungen, in denen andere Systeme die Prüfung anstoßen sollen: eine Pipeline, die nach jeder Änderung am Data Product einen Freigabebericht erzeugt, ein Controlling-Werkzeug, das den erklärten Anteil einer Antwort anzeigt, oder ein KI-Agent, der vor einer Handlung die Freigabehindernisse seiner Stufe abfragt. Dieser Abschnitt beschreibt die Erweiterung, die das Prüfverfahren als zustandslose HTTP-Schnittstelle bereitstellt. Sie ist mit dem Web-Framework FastAPI umgesetzt und liegt im Verzeichnis \texttt{web/} neben \texttt{src/}. Die Bausteine aus \texttt{src/} werden dabei nicht verändert; die Schnittstelle ruft ausschließlich deren öffentliche Funktionen auf.

Der Abschnitt ist wie folgt aufgebaut. Zunächst werden die Entwurfsziele und die Architektur dargestellt, dann die Endpunkte und die Datenverträge. Danach werden die formalen Eigenschaften der Schnittstelle herausgearbeitet, die sich aus den Sätzen des Abschnitts~\ref{sec:operationalization} ergeben, und das Fehlermodell sowie die Maßnahmen zur Robustheit erläutert. Ein Beispiel führt die Fallstudie über die Schnittstelle durch. Abschließend werden Qualitätssicherung, Einsatzszenarien und Grenzen behandelt.

\subsection{Entwurfsziele}
Die Schnittstelle folgt den Leitplanken der Arbeit und übersetzt sie in Anforderungen an einen Netzdienst. Fünf Ziele bestimmen den Entwurf.

\paragraph{Keine zusätzliche Fachlogik.}
Jede Rechnung geschieht in \texttt{src/}. Die Schnittstelle liest Eingaben, ruft eine Bibliotheksfunktion auf und serialisiert das Ergebnis. Dadurch kann es keine zweite, abweichende Implementierung der Freigabelogik geben, und Änderungen an den Definitionen der Arbeit wirken automatisch auch auf die Schnittstelle. Das ist zugleich das Gegenmittel gegen die Stolperfalle, dass verschiedene Werkzeuge dieselbe Kennzahl unterschiedlich berechnen (Abschnitt~\ref{sec:pitfalls}, Nr.~4).

\paragraph{Zustandslosigkeit und Reproduzierbarkeit.}
Jede Anfrage enthält ihre Eingaben vollständig, auch den Stichtag. Der Server hält keine Sitzung, speichert keine Spezifikation und keine Antwort und liest nicht die Systemuhr. Damit überträgt sich der Reproduzierbarkeitssatz der Bibliothek (Satz~\ref{thm:repro}) auf die Schnittstelle.

\paragraph{Exaktheit.}
Alle Rechnungen verwenden exakte Brüche. Zahlen kommen als JSON-Zahl oder als Text (\texttt{"3/20"}, \texttt{"0.15"}) und gehen als Bruch-Text hinaus. Es gibt keine Gleitkommarundung, die einen Schwellenvergleich wie $\bar L\le\varepsilon W$ am Rand verfälschen könnte.

\paragraph{Strenge Eingabe.}
Unbekannte Felder, falsche Typen und ungültige Zahlen werden abgewiesen, nichts wird stillschweigend ersetzt. Das entspricht dem Grundsatz, dass ein fehlender oder unklarer Nachweis nie als Erfolg zählt.

\paragraph{Dokumentierte, prüfbare Verträge.}
Jeder Endpunkt ist in einem maschinenlesbaren Schema (OpenAPI) beschrieben, mit Beispielen und Fehlerantworten. Das Schema ist selbst Gegenstand von Tests, damit Dokumentation und Verhalten nicht auseinanderlaufen.

\subsection{Architektur}
Die Schnittstelle besteht aus drei Modulen in einem eigenen Paket \texttt{web}, die auf das Paket \texttt{aidataeco} (das Verzeichnis \texttt{src/}) aufsetzen. Tabelle~\ref{tab:api-layers} zeigt die Schichten.

\begin{table}[htbp]
\centering
\small
\caption{Schichten der JSON-API}
\label{tab:api-layers}
\begin{tabular}{P{0.24\linewidth}P{0.66\linewidth}}
\toprule
\textbf{Schicht} & \textbf{Aufgabe} \\
\midrule
\texttt{src/} (Paket \texttt{aidataeco}) & Fachlogik: Datennetz, Data Product, Error Budget, Diagnose, Governance, Auswertung, Freigabebericht, Vergleich. Unverändert. \\
\texttt{web/convert.py} & Eingabeprüfung und Umwandlung zwischen JSON und den exakten Bausteinen: Zahlen als Brüche, Stufen aus Nummer oder Name, Schutz vor übergroßen Exponenten. \\
\texttt{web/schemas.py} & Request- und Response-Modelle mit Feldbeschreibungen und Beispielen; sie erzeugen das OpenAPI-Schema und validieren die Typen. \\
\texttt{web/api.py} & FastAPI-Anwendung: Endpunkte, Abbildung von Fehlern auf HTTP-Statuscodes, Größenbegrenzung, Startfunktion. \\
\bottomrule
\end{tabular}
\end{table}

Eine Anfrage durchläuft vier Stufen. Erstens validiert FastAPI mit den Modellen aus \texttt{schemas.py} die Typen des Requests. Zweitens prüft \texttt{convert.py} die freien Dokumente (Spezifikation, Antworten, Berichte) auf übergroße Zahlentexte. Drittens liest die Bibliothek die Dokumente streng und führt die Rechnung aus. Viertens wird das Ergebnis in das Antwortmodell serialisiert, das FastAPI erneut validiert. Fehler in jeder dieser Stufen werden in einheitliche Fehlerantworten übersetzt (Abschnitt~\ref{sec:api-errors}).

Die Bibliothek wird unter ihrem Paketnamen \texttt{aidataeco} importiert, so als wäre sie installiert. Das Projekt weist \texttt{src} dazu in der Paketkonfiguration dem Namen \texttt{aidataeco} zu, weil \texttt{src} als Paketname mehrdeutig ist. Die Konfiguration nimmt außerdem das Paket \texttt{web}, die Abhängigkeiten \texttt{fastapi} und \texttt{uvicorn} sowie ein Konsolenskript zum Start des Servers auf.

\subsection{Ressourcen und Endpunkte}
Die Schnittstelle stellt 28 Operationen unter dem Präfix \texttt{/api/v1} bereit. Sie sind nach den Abschnitten der Arbeit gruppiert, auf deren Definitionen sie beruhen. Tabelle~\ref{tab:api-endpoints} gibt den Überblick.

\begin{table}[htbp]
\centering
\small
\caption{Endpunkte der JSON-API und ihr Bezug zur Arbeit}
\label{tab:api-endpoints}
\begin{tabular}{P{0.17\linewidth}P{0.42\linewidth}P{0.31\linewidth}}
\toprule
\textbf{Gruppe} & \textbf{Endpunkte (unter \texttt{/api/v1})} & \textbf{Bezug zur Arbeit} \\
\midrule
\texttt{meta} & \texttt{GET /health}, \texttt{GET /options} & Stufen, Nachweise, Grenzen \\
\texttt{specs} & \texttt{POST /specs/validate} & Spezifikation $\Sigma$ (Def.~\ref{def:spec}) \\
\texttt{reports} & \texttt{POST /reports}, \texttt{/reports/markdown}, \texttt{/reports/cases}, \texttt{/reports/compare}, \texttt{/reports/compare/markdown} & Evidenzvektor, Freigabehindernisse, Regressionsfreiheit (Def.~\ref{def:evidence}, \ref{def:blockers}, \ref{def:regression}) \\
\texttt{evaluation} & \texttt{POST /evaluation/pass-rate}, \texttt{/explained-share}, \texttt{/completeness}, \texttt{/freshness} & Kennzahlen, Abschnitt~\ref{sec:metrics} \\
\texttt{error-budget} & \texttt{POST /error-budget/check}, \texttt{/linear-loss}, \texttt{/decomposition}, \texttt{/missing-bias}, \texttt{/lipschitz-bound} & Gut-genug-Bedingung, Fehlerzerlegung, Abschnitt~\ref{sec:error} \\
\texttt{diagnosis} & \texttt{POST /diagnosis/posterior}, \texttt{/bayes-update}, \texttt{/voi}, \texttt{/deviations}; \texttt{GET /diagnosis/causes} & Bayes-Indikator, Lern-Loop, VOI, Abschnitt~\ref{sec:opportunity} \\
\texttt{governance} & \texttt{GET /governance/levels}; \texttt{POST /governance/blockers} & Vertrauensstufen, Tabelle~\ref{tab:reliability-levels} \\
\texttt{integration} & \texttt{POST /integration/select} & Datenquellen unter Budget (Erweiterung der VOI-Bedingung) \\
\texttt{experiments} & \texttt{POST /experiments/fingerprint}; \texttt{GET /experiments/versions} & Reproduzierbares Protokoll \\
\texttt{case-study} & \texttt{GET /case-study} & Fallstudie, Abschnitt~\ref{sec:case-study} \\
\bottomrule
\end{tabular}
\end{table}

Der Kern ist die Gruppe \texttt{reports}. Der Endpunkt \texttt{POST /reports} nimmt Spezifikation, Antworten, Stichtag, erwarteten Verlust und optional die Versionen entgegen und liefert den Freigabebericht im selben JSON-Format wie die Kommandozeile. \texttt{/reports/markdown} liefert dieselbe Auswertung als lesbaren Text. \texttt{/reports/cases} wertet jede Referenzfrage einzeln aus, also Antwort, Bestehen innerhalb der Toleranz und Vollständigkeit der Provenance, und eignet sich als Diagnosehilfe, bevor ein Bericht gebildet wird. \texttt{/reports/compare} vergleicht zwei gespeicherte Berichte und meldet, ob eine Regression im Sinn von Definition~\ref{def:regression} vorliegt.

Die übrigen Gruppen machen einzelne Rechenschritte der Arbeit separat zugänglich. Das ist aus zwei Gründen sinnvoll. Erstens verwenden Teams häufig nur einen Teil des Verfahrens, etwa die Gut-genug-Bedingung oder die gewichtete Bestehensquote, und wollen dafür keine vollständige Spezifikation aufbauen. Zweitens lassen sich die Formeln der Arbeit so nachrechnen: Wer den erklärten Anteil der Fallstudie prüfen will, sendet die drei Ursachen an \texttt{/evaluation/explained-share} und erhält $3/22\approx 13{,}6\,\%$.

Zwei Endpunkte verdienen eine Bemerkung. \texttt{GET /governance/levels} leitet die Anforderungen jeder Stufe nicht aus einer zweiten Tabelle ab, sondern ruft die Freigabelogik mit vollständig besetzten Rollen und ohne jeden Nachweis auf; die Hindernisse sind dann genau die Anforderungen der Stufe. Eine Änderung der Anforderungen in der Bibliothek erscheint damit ohne weiteren Eingriff in der Schnittstelle. \texttt{POST /integration/select} löst das Auswahlproblem für Datenquellen mit dem Entscheider des Referenzprojekts; wegen des Rechenaufwands ist die Optionszahl begrenzt (Abschnitt~\ref{sec:api-robust}).

\subsection{Datenverträge}
\paragraph{Zahlen.}
Eine Zahl ist eine JSON-Zahl oder ein Text in Dezimal- oder Bruchschreibweise. Sie wird in einen exakten Bruch $q\in\QQ$ umgewandelt; Gleitkommawerte werden über ihre Dezimaldarstellung gelesen (\texttt{0.1} wird zu $1/10$). Ergebnisse stehen als Text, ganze Zahlen ohne Nenner (\texttt{"3"}), sonst als Bruch (\texttt{"3/22"}). Wahrheitswerte sind keine Zahlen. Diese Konvention stimmt mit der Darstellung der Berichte in der Kommandozeile überein, sodass ein gespeicherter Bericht unverändert an \texttt{/reports/compare} übergeben werden kann.

\paragraph{Datum und Stichtag.}
Daten sind ISO-Texte der Form \texttt{JJJJ-MM-TT}. Der Stichtag $\tau$ ist ein Pflichtfeld von \texttt{POST /reports}; es gibt bewusst keinen Standardwert, weil ein Standardwert die Systemuhr einführen würde.

\paragraph{Dokumente.}
Spezifikation, Antworten und Berichte sind frei strukturierte JSON-Objekte mit denselben Schlüsseln wie die YAML- oder JSON-Dateien der Kommandozeile. Ihre Typen sind im Schema nur als Objekt beschrieben, ihr Inhalt wird von der Bibliothek geprüft. Das vermeidet, die Struktur zweimal zu pflegen, und garantiert, dass die Schnittstelle genau die Dokumente akzeptiert, die auch die Kommandozeile akzeptiert. Die Beispiele im OpenAPI-Schema zeigen die Struktur.

\paragraph{Stufen.}
Eine Stufe wird als Nummer (1 bis 4) oder als Name (\texttt{HINT}, \texttt{CROSS\_CHECK}, \texttt{DECISION\_PREPARATION}, \texttt{AUTOMATED\_ACTION}) angegeben. \texttt{GET /options} liefert beide Darstellungen sowie die Menge der messbaren Nachweise, die in einer Spezifikation nicht erklärt werden dürfen, und der organisatorischen Nachweise, die erklärt werden dürfen.

\subsection{Formale Eigenschaften}
\label{sec:api-formal}
Die Schnittstelle fügt dem Verfahren keine neue Mathematik hinzu, aber sie erlaubt es, Eigenschaften der Bibliothek auf die Netzschicht zu übertragen. Sei $\mathrm{Anf}$ die Menge der wohlgeformten Anfragen eines Endpunkts $e$ und $\mathrm{Ant}_e$ die Menge seiner Erfolgsantworten. Ein Handler ist eine Abbildung
\[
H_e:\ \mathrm{Anf}\ \to\ \mathrm{Ant}_e\ \cup\ \mathrm{Fehler}.
\]

\begin{satz}[Treue des Berichts-Endpunkts]\label{thm:api-fidelity}
Für jede Anfrage $r=(\Sigma,a,\tau,\bar L,v)$, für die $H_{\mathrm{reports}}$ eine Erfolgsantwort liefert, ist diese Antwort gleich der JSON-Darstellung des Freigabeberichts, den die Bibliothek zu denselben Eingaben bildet. Insbesondere stimmen Freigabehindernisse $B(s)$, Evidenzvektor $e$ und Fingerprint überein, und die Markdown-Antwort ist der Markdown-Text desselben Berichts.
\end{satz}
\begin{proof}
Der Handler liest $\Sigma$ und $a$ mit den Lesefunktionen der Bibliothek, wandelt $\tau$ und $\bar L$ in ein Datum beziehungsweise einen exakten Bruch um und ruft die Berichtsfunktion auf. Das Ergebnis wird mit der Serialisierungsfunktion der Bibliothek ausgegeben; die Antwortmodelle ändern keine Werte, sondern beschreiben und validieren sie. Der Kommandozeilenaufruf verwendet dieselben Funktionen. Damit sind beide Ausgaben für gleiche Eingaben gleich.
\end{proof}

Der Satz ist eine Eigenschaft der Implementierung und kein Satz über Datennetze; er ist deshalb nicht bewiesen im Sinn der Formeln, sondern durch Konstruktion begründet und durch einen Test abgesichert, der Antwort und Bibliotheksergebnis vergleicht (Abschnitt~\ref{sec:api-tests}).

\begin{satz}[Zustandslosigkeit und Reproduzierbarkeit der Antwort]\label{thm:api-stateless}
Zwei Anfragen mit gleichem Inhalt, einschließlich der angegebenen Versionen $v$, liefern dieselbe Antwort, unabhängig von der Reihenfolge und der Zahl früherer Anfragen. Sind die Versionen nicht angegeben, gilt dies, solange sich die installierten Versionen von Python und Abhängigkeiten auf dem Server nicht ändern.
\end{satz}
\begin{proof}
Jeder Handler ist eine Funktion seiner Anfrage: Er liest keine Datei, speichert nichts und verwendet weder Zufall noch Systemzeit; der Stichtag ist Eingabe. Die Bibliotheksfunktionen sind nach Satz~\ref{thm:repro} Funktionen ihrer Eingaben. Ohne Angabe von $v$ wird $v$ durch die installierten Versionen ersetzt, die damit zur Eingabe werden.
\end{proof}

Aus Satz~\ref{thm:api-stateless} folgt, dass sich ein Bericht nachträglich bestätigen lässt: Wer die Anfrage aufbewahrt, kann an beliebiger Stelle denselben Bericht und denselben Fingerprint erzeugen. Für diese Eigenschaft sollten Clients die Versionen stets angeben.

\begin{satz}[Unbekanntes zählt auch an der Schnittstelle nie als Erfolg]\label{thm:api-conservative}
Ein JSON-Wert \texttt{null} oder ein fehlendes optionales Feld wird als $\bot$ gelesen. Eine nicht angegebene Antwort $\hat y_j$ gilt als unbeantwortet, ein nicht angegebener erwarteter Verlust $\bar L$ als nicht belegt, eine nicht angegebene Qualität $q_v$ oder Aktualisierung $t_v$ als unbekannt. Für die Eingaben, die sich so aus einer Anfrage durch Weglassen ergeben, gilt die Monotonie aus Satz~\ref{thm:conservative}: Das Weglassen kann die Menge $B$ der Freigabehindernisse nur vergrößern.
\end{satz}
\begin{proof}
Die Lesefunktionen bilden fehlende und \texttt{null}-Angaben auf $\bot$ ab und ersetzen sie durch keinen Standardwert; das gilt für Antworten, Qualität, Aktualisierung und erwarteten Verlust. Die Ersetzungen (a) bis (e) des Satzes~\ref{thm:conservative} sind genau diese Abbildungen auf $\bot$. Der Satz liefert die Behauptung.
\end{proof}

Ausgenommen sind die organisatorischen Nachweise $o$: Fehlt ein Feld unter \texttt{evidence}, gilt es als nicht erklärt ($0$), eine Erklärung ist aber nicht überprüfbar. Auch das ist unverändert gegenüber der Bibliothek.

\begin{bemerkung}[Kein Teilergebnis]
Für jede Anfrage liefert ein Endpunkt entweder eine vollständige Erfolgsantwort oder einen Fehler; ein Bericht mit nur teilweise ausgewerteten Eingaben wird nicht ausgegeben. Das setzt voraus, dass die Bibliothek ungültige Eingaben ausnahmslos durch einen Eingabefehler meldet. Das ist ihre dokumentierte Vertragsregel; andere Ausnahmen gelten als Programmfehler und führen zu einem Serverfehler (HTTP~500), nicht zu einer scheinbar gültigen Antwort.
\end{bemerkung}

\subsection{Fehlermodell}
\label{sec:api-errors}
Die Schnittstelle unterscheidet fachliche Eingabefehler von Fehlern der Anfrage und von Ausfällen. Tabelle~\ref{tab:api-errors} fasst die Statuscodes zusammen.

\begin{table}[htbp]
\centering
\small
\caption{Statuscodes und Fehlerschemata der JSON-API}
\label{tab:api-errors}
\begin{tabular}{P{0.09\linewidth}P{0.27\linewidth}P{0.50\linewidth}}
\toprule
\textbf{Status} & \textbf{Schema} & \textbf{Bedeutung} \\
\midrule
200 & Antwortmodell des Endpunkts & Erfolg; bei Berichten auch dann, wenn Freigabehindernisse bestehen. Das Hindernis ist Ergebnis, kein Fehler. \\
422 & \texttt{error.\allowbreak type}\,=\,\texttt{input\_error} & Fachlich ungültige Eingabe: unbekannte Felder, falsche Typen im Dokument, ungültige Zahlen oder Daten, doppelte Referenzfragen, erklärte messbare Nachweise, übergroße Exponenten. \\
422 & \texttt{\{"detail":[\ldots]\}} (FastAPI) & Request-Feld fehlt oder hat den falschen Typ, etwa ein Stichtag, der nicht dem Muster \texttt{JJJJ-MM-TT} entspricht. \\
413 & \texttt{error.\allowbreak type}\,=\,\texttt{payload\_too\_large} & \texttt{Content-Length} überschreitet die Grenze. \\
400 & \texttt{error.\allowbreak type}\,=\,\texttt{invalid\_header} & \texttt{Content-Length} ist keine ganze Zahl. \\
503 & \texttt{error.\allowbreak type}\,=\,\texttt{backend\_error} & Der Entscheider des Solvers konnte keine verlässliche Antwort liefern. \\
\bottomrule
\end{tabular}
\end{table}

Die Trennung von Ergebnis und Fehler ist fachlich wichtig. Ein Freigabebericht mit sieben Freigabehindernissen ist eine erfolgreiche Antwort mit dem Statuscode 200: Die Schnittstelle hat ihre Aufgabe erfüllt, indem sie festgestellt hat, dass das Datennetz für diesen Anwendungsfall nicht ausreicht. Würde man ein Hindernis als HTTP-Fehler melden, vermischte man den Zustand des Datennetzes mit dem Zustand des Dienstes, und Überwachungswerkzeuge würden einen Befund als Störung zählen. Clients, die eine Pipeline anhalten wollen, werten das Feld \texttt{blocked} beziehungsweise \texttt{has\_regression} aus, ähnlich wie die Kommandozeile mit dem Exit-Code 1 signalisiert.

Die Fehlermeldungen nennen den Pfad zur fehlerhaften Stelle (etwa \texttt{use\_case.level}) und bleiben deutsch wie die Bibliothek. Sie enthalten keine Stapelüberwachungen und keine Inhalte der Anfrage über das hinaus, was zur Korrektur nötig ist.

\subsection{Robustheit und Sicherheit}
\label{sec:api-robust}
Eine Schnittstelle, die frei strukturierte Dokumente annimmt und exakt rechnet, hat Angriffsflächen, die die Bibliothek allein nicht hat. Die Erweiterung adressiert drei davon und benennt, was sie bewusst nicht leistet.

\paragraph{Übergroße Zahlentexte.}
Die exakte Arithmetik wandelt Texte wie \texttt{"1e999999999"} in einen Bruch mit einer Milliarde Ziffern um und würde Speicher und Rechenzeit verbrauchen. Die Schnittstelle weist Zahlentexte ab, deren Exponent mehr als drei Ziffern hat; die Prüfung läuft über alle Texte der Spezifikation, der Antworten und der Berichte, iterativ und damit unabhängig von der Verschachtelungstiefe. Sie greift nur bei Texten, die insgesamt eine Zahl mit Exponent sind. Zeichenfolgen, die zufällig ein \emph{e} und Ziffern enthalten, etwa ein SHA-256-Fingerprint, bleiben unberührt. Ein erster Entwurf der Prüfung hat solche Fingerprints fälschlich abgewiesen; der Fehler wurde durch den Test des Berichtsvergleichs entdeckt und ist ein Beispiel dafür, dass der Eingabeschutz selbst getestet werden muss.

\paragraph{Übergroße Anfragen.}
Anfragen, deren \texttt{Content-Length} die Grenze überschreitet (Standard 1~MiB, einstellbar über eine Umgebungsvariable), werden mit HTTP~413 abgewiesen, bevor der Inhalt gelesen wird. Übertragungen ohne \texttt{Content-Length} (chunked) erfasst diese Schicht nicht; dafür ist ein vorgeschalteter Reverse Proxy zuständig.

\paragraph{Rechenaufwand des Solvers.}
Der Endpunkt \texttt{POST /integration/select} kodiert das Auswahlproblem für den Entscheider des Referenzprojekts. Dessen Laufzeit wächst stark mit der Zahl der Optionen: In einer Einzelmessung in der Entwicklungsumgebung dauerte die Lösung bei acht Optionen etwa eine Sekunde, bei zehn rund zwei Sekunden und bei zwölf bereits etwa 24 Sekunden. Die Schnittstelle begrenzt daher die Zahl der Optionen auf zehn und Nutzen, Kosten und Budget auf 100\,000 ganzzahlige Einheiten. Die Zahlen sind eine Beobachtung unter bestimmten Bedingungen und kein Laufzeitbeweis; sie begründen die Grenze, nicht eine Garantie. Der Endpunkt ist synchron implementiert und läuft damit in einem Arbeits-Thread des Servers, sodass andere Anfragen nicht blockieren, aber die Auslastung der Arbeits-Threads durch viele gleichzeitige schwere Anfragen bleibt ein Betriebsthema (Ausblick).

\paragraph{Was die Schnittstelle nicht leistet.}
Sie hat keine Authentifizierung, keine Autorisierung, keine Begrenzung der Anfragerate und keine Transportverschlüsselung. Diese Funktionen gehören in die Betriebsumgebung (Gateway, Reverse Proxy, Identitätsanbieter), nicht in eine Rechenschnittstelle, die möglichst wenig Zustand haben soll. Spezifikationen und Antworten können sensible Unternehmensdaten enthalten; die Zustandslosigkeit verringert das Risiko, weil nichts gespeichert wird, ersetzt aber keinen Schutz auf dem Transportweg und in Protokollen des Betriebs. Die Schnittstelle ist deshalb für den Einsatz im geschützten Netz oder hinter einem Gateway gedacht, nicht für den ungeschützten Betrieb im Internet.

\subsection{Beispiel: Die Fallstudie über die Schnittstelle}
\label{sec:api-example}
Das Data Product „Sales Result 360“ der Fallstudie (Abschnitt~\ref{sec:case-study}) dient als durchgehendes Beispiel. Die Spezifikation ist die Beispieldatei des Projekts, die Antworten sind die der Fallstudie: Zur Frage nach der Expressfracht antwortet die KI mit 0 und vollständiger Provenance, die Frage nach dem Produktmix bleibt unbeantwortet, die Rabattfrage wird richtig beantwortet.

\paragraph{Erklärter Anteil.}
Eine Anfrage an \texttt{POST /evaluation/explained-share} mit den drei Ursachen der Abweichung (Expressfracht 230, Produktmix 150, Rabatte 60 mit erklärten Anteilen 0, 0 und 1) liefert:
\begin{lstlisting}
{"explained_share": "3/22", "unexplained_share": "19/22",
 "explained_amount": "60", "total_effect": "440"}
\end{lstlisting}
Das sind die rund 14\,\% der Fallstudie. Dieselbe Zahl liefert \texttt{GET /case-study} mit der Aufteilung der Ursachen.

\paragraph{Gut-genug-Bedingung.}
Mit Stakes $W=440$ und $\varepsilon=1/20$ beträgt das zulässige Budget 22. Der erwartete Verlust $\bar L=380$ aus Abschnitt~\ref{sec:operationalization} überschreitet es deutlich; \texttt{POST /error-budget/check} antwortet:
\begin{lstlisting}
{"limit": "22", "expected_loss": "380",
 "good_enough": false, "headroom": "-358"}
\end{lstlisting}

\paragraph{Freigabebericht.}
Die Anfrage an \texttt{POST /reports} enthält die Spezifikation, die Antworten, den Stichtag 22.~Oktober 2026 und den erwarteten Verlust 380. Die Antwort hat den Statuscode 200 und, in Auszügen, diesen Inhalt:
\begin{lstlisting}
{"product": "Sales Result 360", "level": "DECISION_PREPARATION",
 "missing_links": [["orders", "freight"]],
 "evaluation": {"pass_rate": "1/6",
                "provenance_coverage": "2/3", "unanswered": 1},
 "min_pass_rate": "3/4", "budget_limit": "22",
 "blocked": true, "blockers": [ ... 7 Eintraege ... ],
 "fingerprint": "8b2dbd0b..."}
\end{lstlisting}
Die sieben Freigabehindernisse sind genau die der Fallstudie in Abschnitt~\ref{sec:operationalization}: Die Kennzahlen sind nicht verknüpft, weil die Verbindung zwischen Auftrag und Fracht fehlt; die Provenance ist unvollständig, weil die Antwort zum Produktmix fehlt; die Qualität des Knotens \texttt{products} ist nicht messbar; das Error Budget ist überschritten; die Knoten \texttt{freight} und \texttt{products} sind nicht aktuell zum Stichtag; und die Bestehensquote $1/6$ liegt unter dem Mindestwert $3/4$. Die Schnittstelle bestätigt damit das Ergebnis der Bibliothek auf der Netzebene.

\paragraph{Regressionsvergleich.}
Nach dem Anbinden der Frachtdaten und dem Ergänzen der Produktgruppe wird ein zweiter Bericht gebildet. Beide Berichte werden an \texttt{POST /reports/compare} übergeben; die Antwort nennt je Kennzahl die Richtung, die behobenen und die neu hinzugekommenen Hindernisse und das Feld \texttt{has\_regression}. Eine Pipeline kann daran einen Merge oder eine Freigabe knüpfen. Wegen der Zustandslosigkeit speichert der Aufrufer die Berichte selbst, etwa als Artefakte seines Build-Systems.

\paragraph{Fehlerfall.}
Wird als erwarteter Verlust \texttt{"1e5000"} übergeben, antwortet die Schnittstelle mit HTTP~422 und
\begin{lstlisting}
{"error": {"type": "input_error", "message":
  "expected_loss: Exponenten mit mehr als 3 Ziffern sind nicht zulaessig"}}
\end{lstlisting}
und rechnet nicht.

\subsection{Qualitätssicherung}
\label{sec:api-tests}
Die Schnittstelle wird wie die Bibliothek mit Tests abgesichert. Die Testsuite des Projekts umfasst 291 Tests; davon prüfen 106 die Schnittstelle, 185 die Bausteine, die Kommandozeile und die Beispiele der Dokumentation. Zeilen- und Zweigabdeckung betragen für \texttt{src/} und für \texttt{web/} jeweils 100\,\%. Vier Testarten sind zu unterscheiden:
\begin{itemize}
\item \textbf{Treue-Tests} vergleichen die Antworten mit der Bibliothek und der Kommandozeile: Der Bericht der Schnittstelle ist gleich dem JSON der Kommandozeile, die Markdown-Antworten gleich den Markdown-Texten, die Blocker einer Stufe gleich den Ergebnissen der Freigabelogik. Sie stützen Satz~\ref{thm:api-fidelity}.
\item \textbf{Verhaltenstests} prüfen die erwarteten Werte der Arbeit, etwa den erklärten Anteil $3/22$ der Fallstudie, die Grenze der Gut-genug-Bedingung, die Posterior-Odds des Bayes-Indikators und die Zahl der Anforderungen je Stufe.
\item \textbf{Fehlertests} decken jeden Fehlerzweig ab: ungültige Daten, unbekannte Fragen, erklärte messbare Nachweise, doppelte Fragen, übergroße Exponenten, Größenlimit, ungültigen Header und den simulierten Ausfall des Solvers.
\item \textbf{Vertragstests} prüfen das OpenAPI-Schema: Jede Operation hat Zusammenfassung, Beschreibung, Tag und Fehlerantworten, jedes Feld eine Beschreibung, jeder Tag ist deklariert und benutzt; die Namen der Stufen und Ursachen im Schema stimmen mit den Aufzählungen der Bibliothek überein, und die Version der Schnittstelle gleicht der Projektversion.
\end{itemize}
Satz~\ref{thm:api-stateless} wird durch einen Test gestützt, der dieselbe Anfrage vor und nach anderen Anfragen wiederholt und gleiche Antworten verlangt, sowie durch Tests des Fingerprints, der sich bei geänderten Eingaben oder Versionen ändern muss.

Die Abdeckung ist ein Hinweis auf die Vollständigkeit der Prüfung, kein Nachweis der Korrektheit. Ein Zweig kann ausgeführt werden, ohne dass sein Ergebnis fachlich geprüft wird; deshalb vergleichen die Treue- und Verhaltenstests Werte und nicht nur Statuscodes. Ebenso beweisen die Tests nicht, dass die Freigabelogik die Wirklichkeit eines Unternehmens abbildet; sie belegen nur, dass die Schnittstelle die Logik der Arbeit treu wiedergibt.

\subsection{Einsatzszenarien}
\label{sec:api-use}
Die folgenden Szenarien zeigen, wie die Schnittstelle in die Abläufe der Arbeit passt. Sie sind Entwurfsvorschläge, keine in der Praxis erprobten Verfahren.

\paragraph{Freigabe-Gate in der Datenpipeline.}
Nach jeder Änderung am Data Product (neue Quelle, geänderte Definition, neue Version) erzeugt die Pipeline einen Bericht mit \texttt{POST /reports} und vergleicht ihn mit dem letzten gespeicherten Bericht über \texttt{POST /reports/compare}. Eine Regression hält die Änderung an. Das setzt den Gedanken um, dass Regelmäßigkeit Teil der Verantwortung des Data Product Owners ist (Abschnitt~\ref{sec:release}).

\paragraph{Anzeige in Controlling-Werkzeugen.}
Ein Dashboard fragt \texttt{/reports/cases} ab und zeigt je Referenzfrage Bestehen und Provenance. Die Beteiligten sehen, welche Fragen die KI nicht beantwortet, bevor sie ihrer Antwort auf eine Geschäftsfrage vertrauen. Das entspricht der Forderung, Definition, Quelle und Datenstichtag einzufordern (Abschnitt~\ref{sec:pitfalls}, Nr.~1).

\paragraph{Vorprüfung durch KI-Agenten.}
Ein Agent, der eine Handlung auslösen soll, ruft \texttt{POST /governance/blockers} mit der Stufe seines Anwendungsfalls auf und führt die Handlung nur aus, wenn die Liste leer ist und ein Mensch die Freigabe erklärt hat. Die Schnittstelle ersetzt dabei nicht die Entscheidung des Use-Case Owners (Abschnitt~\ref{sec:api-formal}); sie stellt nur sicher, dass der Agent die formalisierten Bedingungen kennt und nicht umgeht. Der Aufruf könnte als Werkzeug eines Agenten bereitgestellt werden, etwa über das Model Context Protocol; eine solche Anbindung ist nicht Teil der Umsetzung.

\paragraph{Planung von Datenintegrationen.}
Die Gruppen \texttt{diagnosis} und \texttt{integration} unterstützen die ökonomische Priorisierung von Datenverbesserungen: Nach einer Abweichung berechnet \texttt{/diagnosis/voi}, ob die Integration einer Quelle ihre Kosten rechtfertigt, und \texttt{/integration/select} wählt aus mehreren Kandidaten die Auswahl mit dem größten Nettonutzen unter einem gemeinsamen Budget.

\subsection{Grenzen der Schnittstelle}
\label{sec:api-limits}
Die Schnittstelle übernimmt alle Grenzen der Bibliothek: Erklärte organisatorische Nachweise werden nicht überprüft, die Provenance ist nur formal, nicht inhaltlich geprüft, der erwartete Verlust wird geliefert und nicht geschätzt, der Vergleich kennt kein Messrauschen, und ein Bericht ohne Freigabehindernisse ist keine Freigabe. Hinzu kommen Grenzen, die aus der Netzschicht folgen.
\begin{itemize}
\item Die Schnittstelle speichert nichts. Historie, Versionierung von Berichten und ein Prüfpfad (Audit-Log) liegen beim Aufrufer.
\item Es gibt keine Authentifizierung, Autorisierung und Ratenbegrenzung (Abschnitt~\ref{sec:api-robust}).
\item Anbindungen an konkrete KI-Anwendungen sind nicht enthalten; die Antworten der KI müssen dem Verfahren als Dokument übergeben werden. Die Bibliothek enthält dafür die Schnittstelle des Antwortgebers, die Schnittstelle bietet sie nicht an.
\item Die Größenbegrenzung betrifft nur die angekündigte Länge, und die Rechenzeit der Endpunkte außer der Auswahl von Datenquellen ist nicht eigens begrenzt. Sehr große Spezifikationen werden zwar korrekt, aber langsam verarbeitet.
\item Die Version 1 des Pfades \texttt{/api/v1} verspricht nicht, dass spätere Versionen der Arbeit die Berichtsstruktur nicht ändern; für gespeicherte Berichte ist daher der Fingerprint samt Versionen aufzubewahren.
\end{itemize}

\FloatBarrier

\section{Checkliste, Diskussion, Schlussfolgerung}

\subsection{Checkliste: Ist diese KI-Antwort vertrauenswürdig?}
\label{sec:checklist}

Vor der Verwendung einer KI-Antwort sollten mindestens die folgenden Punkte geprüft und für den konkreten Einsatzfall beantwortet werden:
\begin{enumerate}[label=$\square$]
\item Ist die Frage klar formuliert, und sind Stakes sowie möglicher Schaden einer Fehlentscheidung bekannt?
\item Sind die wesentlichen Einflussgrößen im Datennetz vorhanden und miteinander verknüpft?
\item Ist die verwendete Kennzahl eindeutig definiert, und ist ein Data Owner benannt?
\item Nennt die Antwort Definition, Datenquelle und Datenstichtag?
\item Sind Qualität und Aktualität der entscheidenden Daten bekannt und für den Anwendungsfall ausreichend?
\item Wurde die Antwort mit der eigenen Analyse oder Strategie gegengeprüft?
\item Ist der mögliche verbleibende Fehler im Verhältnis zu den Folgen der Entscheidung vertretbar?
\item Ist festgelegt, wann ein Mensch entscheidet, genehmigt oder eingreifen muss?
\end{enumerate}
Ein fehlender Nachweis ist selbst ein relevantes Ergebnis: Er bedeutet nicht automatisch, dass die Antwort falsch ist, aber dass sie nicht ohne weitere Prüfung für eine höher riskante Verwendung freigegeben werden sollte.

\subsection{Diskussion und Grenzen}\label{sec:discussion}
Die formalen Modelle vereinfachen reale Unternehmensprozesse bewusst: Einflussbeziehungen sind nicht immer linear, Referenzwerte können selbst unsicher sein und Ursachen interagieren. Auch die im Fallbeispiel genannten Beträge sind eine illustrative Webinar-Fallstudie und keine unabhängig erhobene Unternehmensanalyse. Die angegebenen Reife- und Fehlermaße müssen daher als fallbezogene Entscheidungshilfen verstanden werden.

Eine weitere Grenze ist die Auswahl der Testfragen. Ein hoher Messwert für leicht beantwortbare oder häufige Fragen sagt wenig über seltene, aber folgenreiche Ausnahmen aus. Testsets sollten deshalb nach möglichen Auswirkungen geschichtet und auch nach Änderungen an Datenquellen, Definitionen, Modellen oder Berechtigungen erneut ausgeführt werden. Ebenso ist eine Abweichung nicht automatisch ein Beleg gegen das Datennetz: Sie kann auf einen Fehler in der internen Analyse, eine andere zulässige Definition oder einen inzwischen veränderten Datenstand hinweisen.

Die Schnittstelle aus Abschnitt~\ref{sec:api} verändert diese Grenzen nicht, sondern macht sie leichter sichtbar: Sie liefert exakt dieselben Ergebnisse wie die Bibliothek (Satz~\ref{thm:api-fidelity}), prüft aber ebenfalls nur, was formalisiert ist. Eine einfach erreichbare Prüfung birgt zudem das Risiko, dass ein formal bestandener Freigabebericht als inhaltliche Garantie missverstanden wird.

Die wirtschaftliche Logik bleibt dennoch belastbar: Datenverbesserungen sollten dort priorisiert werden, wo ihr erwarteter zusätzlicher Nutzen die Integrations- und Betriebskosten rechtfertigt. In der Praxis sind dafür kontrollierte Experimente, historische Was-wäre-wenn-Analysen und eine nachvollziehbare Messung der Folgen nötig. Die Schätzung dieses Nutzens in realen Geschäftsfällen bleibt eine zentrale offene Aufgabe.

\subsection{Ausblick}\label{sec:outlook}
Der Ausblick gliedert sich in fünf Richtungen, die auf den vorhandenen Bausteinen aufbauen: empirische Validierung, Erweiterung des Modells, Ausbau der Schnittstelle, organisatorische Einbettung und Öffnung für KI-Agenten. Tabelle~\ref{tab:roadmap} ordnet sie nach Voraussetzung und erwartetem Ertrag.

\subsubsection{Empirische Validierung}
Die wichtigste offene Aufgabe ist die Überprüfung an realen Fällen. Nötig sind Fallstudien in Unternehmen unterschiedlicher Größe (vgl.\ Abschnitt~\ref{sec:company-size}), in denen Freigabebericht, Fehlerbudget und Bestehensquote über mehrere Zyklen erhoben und mit tatsächlich eingetretenen Fehlentscheidungen verglichen werden. Erst dann lässt sich beurteilen, ob die Vertrauensstufen Fehler tatsächlich trennen und ob die gewählten Schwellen $\theta$ sinnvoll sind. Ergänzend sollten Testsätze nach Auswirkungen geschichtet und der Einfluss des Testsatzes auf $p$ gesondert gemessen werden, damit eine Änderung des Datennetzes von einer Änderung der Referenzfragen unterscheidbar wird (vgl.\ Bemerkung zu den Grenzen des Vergleichs).

\subsubsection{Erweiterung des Modells}
Das Modell rechnet exakt und kennt keine Toleranz. Drei Erweiterungen liegen nahe: erstens ein statistischer Vergleich, der bei kleinen Testmengen Messrauschen berücksichtigt (etwa Konfidenzintervalle für $p$); zweitens nichtlineare und interagierende Einflussgrößen, für die der Anhang bereits eine Skizze enthält; drittens eine abgestufte Aktualität, bei der Daten nicht nur aktuell oder nicht aktuell sind, sondern ihr Alter in die Fehlerabschätzung eingeht. Auch die Schätzung der mittleren Stakes $\bar L$ sollte methodisch abgesichert werden, da das Fehlerbudget linear von ihr abhängt.

\subsubsection{Ausbau der Schnittstelle}
Die Schnittstelle ist bewusst zustandslos. Für den dauerhaften Betrieb sind folgende Ergänzungen sinnvoll: eine Authentifizierung und Rollenprüfung entsprechend den Verantwortlichkeiten, eine revisionssichere Ablage der Freigabeberichte samt Fingerprint, asynchrone Verarbeitung für den lösungsintensiven Teil, dessen Laufzeit exponentiell wächst und derzeit durch Größenlimits begrenzt wird, sowie Webhooks, die bei einer Regression automatisch eine Meldung auslösen. Eine solche Ablage würde die Zustandslosigkeit nicht aufheben, wenn sie als getrennte Komponente vor oder hinter der Schnittstelle liegt und deren Berechnungen unverändert lässt.

\subsubsection{Organisatorische Einbettung}
Die Berichte entfalten Wirkung erst, wenn sie in bestehende Abläufe eingebunden sind: als Prüfschritt in der Auslieferungskette (Freigabe nur bei bestandener Prüfung), als regelmäßiger Bericht an die benannten Verantwortlichen und als Eingang in Governance-Gremien. Hier ist offen, wie viel Automatisierung der Freigabe vertretbar ist; bei hohen Stakes bleibt die menschliche Entscheidung maßgeblich, die Schnittstelle liefert dafür den Nachweis, nicht das Urteil.

\subsubsection{Öffnung für KI-Agenten}
Da die Schnittstelle ein maschinenlesbares Schema (OpenAPI) bereitstellt, können KI-Agenten sie als Werkzeug aufrufen, etwa über das Model Context Protocol. Ein Agent könnte vor einer Antwort prüfen, ob die benötigten Nachweise vorliegen, und andernfalls die Antwort als vorläufig kennzeichnen. Das verlagert die Gegenprüfung in den Antwortprozess selbst. Offen bleibt, wie sich verhindern lässt, dass ein Agent die Prüfung lediglich zur Rechtfertigung nutzt; die konservative Eigenschaft (Satz~\ref{thm:api-conservative}) ist dafür eine notwendige, aber keine hinreichende Bedingung.

\begin{table}[H]
\centering
\caption{Ausblick: Richtungen, Voraussetzungen und erwarteter Ertrag}\label{tab:roadmap}
\begin{tabular}{P{3.6cm}P{4.6cm}P{5.8cm}}
\toprule
Richtung & Voraussetzung & Erwarteter Ertrag \\
\midrule
Empirische Validierung & Zugang zu Fallstudien, historische Fehlerdaten & Belastbare Schwellen, Prüfung der Vertrauensstufen \\
Modellerweiterung & Statistische Methoden, Daten zu Aktualität & Robuster Vergleich bei kleinen Testmengen \\
Ausbau der Schnittstelle & Betriebsumgebung, Rollenkonzept & Revisionssichere, automatisierbare Freigabe \\
Organisatorische Einbettung & Benannte Verantwortliche, Governance & Verbindliche Nutzung des Berichts \\
Öffnung für KI-Agenten & Maschinenlesbares Schema, Agenten-Anbindung & Gegenprüfung im Antwortprozess \\
\bottomrule
\end{tabular}
\end{table}

\subsection{Schlussfolgerung}
Die Frage „Wann ist das Datennetz gut genug für KI?“ hat keine allgemeingültige, absolute Schwelle. Gut genug ist es für einen bestimmten Anwendungsfall, wenn relevante Einflussgrößen verknüpft und definiert sind, Qualität und Aktualität bekannt sind, Antworten nachvollziehbare Provenance tragen und der verbleibende Fehler im Verhältnis zu den Stakes vertretbar ist.

Aus der Fallstudie und dem formalen Modell folgen sieben praktische Kernaussagen:
\begin{enumerate}
\item Die Qualität einer KI-Erklärung hängt maßgeblich von den verfügbaren Daten, ihren Verbindungen und ihrer fachlichen Bedeutung ab, nicht allein vom Modell.
\item „Gut genug“ ist eine wirtschaftliche und anwendungsbezogene Bedingung: Je folgenreicher die Entscheidung, desto stärker müssen Fehler begrenzt und belegt werden.
\item Vertrauen lässt sich schrittweise aufbauen – durch klare Definitionen, messbare Qualität, Provenance und benannte Verantwortliche.
\item KI kann bereits heute als Hinweis und Gegenprüfung nützlich sein. Eine Abweichung zeigt mögliche Datenlücken auf, kann aber ebenso eine unvollständige eigene Analyse offenlegen.
\item Ein Data Product schafft eine wiederverwendbare Grundlage für verschiedene Anwendungen; der Aufwand für gemeinsame Daten und Definitionen kann dadurch mehrfach Nutzen stiften.
\item Vor einem Go-live müssen fachliche Definition, technische Lieferung und Freigabe des konkreten Anwendungsfalls jeweils einer verantwortlichen Person zugeordnet sein.
\item Das Prüfverfahren lässt sich als zustandslose, vollständig getestete Schnittstelle bereitstellen, ohne seine Ergebnisse zu verändern; so wird die Prüfung wiederholbar und in Abläufe einbindbar, ersetzt aber nicht die fachliche Verantwortung.
\end{enumerate}

\begin{quote}
\textbf{Leitgedanke:} Nicht nur fragen, ob die KI gut genug ist, sondern ob das Datennetz für diese Frage gut genug ist. Bis dahin die KI als Hinweis nutzen, ihre Antworten mit der eigenen Analyse abgleichen und aus jeder Abweichung lernen.
\end{quote}

\section*{Danksagung}
Die Arbeit basiert auf den Ideen und Folien von ai-data-network-economic-approach-en.md sowie auf den wissenschaftlichen Vorarbeiten des Autors zur wirtschaftlichen Bewertung eines SAT-Solvers und zu Wirtschaftssystemen unter Unsicherheit \cite{sat2026,uncertainty2026}.
\appendix
\section{Glossar}
\begin{description}[style=nextline]
\item[Datennetz] Gesamtheit der Unternehmensdaten und der fachlichen sowie technischen Verbindungen zwischen ihnen.
\item[Data Product] Wiederverwendbare Einheit aus Daten, ihrer Bedeutung und einem Label mit Nutzungsvoraussetzungen und benannter Verantwortung.
\item[DB\,I / DB\,II] Deckungsbeitrag I beziehungsweise II. In der Fallstudie bezeichnet „Marge“ die DB-II-Marge.
\item[Data Owner] Fachlich verantwortliche Person für Bedeutung und Definition von Kennzahlen.
\item[Data Steward] Unterstützt den Data Owner bei der laufenden Pflege und Anwendung fachlicher Regeln.
\item[Data Product Owner] Verantwortet die Lieferung, Vollständigkeit, Aktualität, Prüfung und Berechtigung des Data Products.
\item[Use-Case Owner] Verantwortet Freigabe und Testfragen einer KI-Anwendung, ihr Error Budget und die Grenze menschlicher Entscheidungen.
\item[Granularität] Detailstufe eines Datensatzes, beispielsweise eine Zeile je Auftragsposition.
\item[Provenance] Nachvollziehbare Angaben zu Definition, Quelle, Version und Datenstichtag einer Antwort.
\item[Error Budget] Für einen Anwendungsfall zulässiger Fehler, bestimmt durch Stakes und mögliche Folgen.
\item[Gegenprüfung (Cross-check)] Vergleich einer KI-Antwort mit eigener Analyse, Strategie oder fachlich geprüfter Referenz.
\item[MCP] Model Context Protocol; ein Standard, über den KI-Anwendungen auf Daten und Werkzeuge zugreifen können.
\item[Value of Information (VOI)] Erwarteter zusätzlicher Nutzen einer Information oder Datenverbesserung im Vergleich zu ihren Kosten.
\item[API (Schnittstelle)] Maschinenlesbarer Zugang zum Prüfverfahren über HTTP mit JSON-Ein- und -Ausgaben.
\item[FastAPI] Python-Framework für Web-Schnittstellen; erzeugt aus den Datenverträgen automatisch eine OpenAPI-Beschreibung.
\item[OpenAPI] Standardisierte, maschinenlesbare Beschreibung einer HTTP-Schnittstelle.
\item[Zustandslos] Jede Anfrage enthält alle Eingaben; der Dienst speichert nichts zwischen Anfragen.
\end{description}

\section{Formale Beweise und Erweiterungen}
\subsection{Erweiterung: Nichtlineare Einflussmodelle}
Die lineare Darstellung kann durch eine allgemeine nichtlineare Regressionsdarstellung $y^*=g(S_q)$ mit Lipschitz-Bedingung $(\exists L: |g(x)-g(x')|\le L\|x-x'\|)$ verallgemeinert werden. Dann gilt eine analoge Fehlerabschätzung:
\[
|\hat y-y^*| \le L\cdot \|S_q-\tilde S_q\| + |B_{def}| + |\eta|.
\]

\subsection{Beweisskizze: VOI-Berechnung}
Mit Posterior $\pi_1(\theta)$ ist der erwartete Wert der Integration eines zusätzlichen Datenobjekts $z$:
\[
\E_{\pi_1}[\Delta W] = \sum_{\theta} \pi_1(\theta) \big( W(\text{with }z,\theta)-W(\text{without }z,\theta)\big).
\]
Dies kann numerisch über Simulation oder historische "what-if"-Analysen geschätzt werden.

\bibliographystyle{plain}
\begin{thebibliography}{99}
\bibitem{webinar2026} Michael Sullmann (BE-terna), Webinar "Data Product and AI in Controlling", Okt 2026.
\bibitem{sat2026} Stephan Epp. \emph{Wirtschaftliche Konsequenzen eines effizienten SAT-Solvers}. 12. Oktober 2026. \url{https://github.com/naphets86/subgraph-sat-solver-econ}.
\bibitem{uncertainty2026} Stephan Epp. \emph{Wirtschaftssysteme unter Unsicherheit}. 3. September 2026. \url{https://github.com/naphets86/odd/tree/main/science/uncrtecos}.

\bibitem{abraham2019} Abraham, R., Schneider, J., vom Brocke, J.: Data governance: A conceptual framework, structured review, and research agenda. \emph{International Journal of Information Management} 49, 424--438, 2019.
\bibitem{amodei2016} Amodei, D., Olah, C., Steinhardt, J., Christiano, P., Schulman, J., Mané, D.: Concrete Problems in AI Safety. arXiv:1606.06565, 2016.
\bibitem{batini2009} Batini, C., Cappiello, C., Francalanci, C., Maurino, A.: Methodologies for Data Quality Assessment and Improvement. \emph{ACM Computing Surveys} 41(3), Article 16, 2009.
\bibitem{breck2017} Breck, E., Cai, S., Nielsen, E., Salib, M., Sculley, D.: The ML Test Score: A Rubric for ML Production Readiness and Technical Debt Reduction. In: \emph{Proceedings of the IEEE International Conference on Big Data}, 1123--1132, 2017.
\bibitem{buneman2001} Buneman, P., Khanna, S., Tan, W.-C.: Why and Where: A Characterization of Data Provenance. In: \emph{Database Theory -- ICDT 2001}, LNCS 1973, 316--330, Springer, 2001.
\bibitem{dama2017} DAMA International: \emph{DAMA-DMBOK: Data Management Body of Knowledge}. 2. Auflage, Technics Publications, 2017.
\bibitem{dehghani2022} Dehghani, Z.: \emph{Data Mesh: Delivering Data-Driven Value at Scale}. O'Reilly Media, 2022.
\bibitem{euaiact2024} Europäische Union: Verordnung (EU) 2024/1689 des Europäischen Parlaments und des Rates vom 13. Juni 2024 zur Festlegung harmonisierter Vorschriften für künstliche Intelligenz (KI-Verordnung). \emph{Amtsblatt der Europäischen Union}, L, 12.7.2024. \url{http://data.europa.eu/eli/reg/2024/1689/oj}.
\bibitem{gebru2021} Gebru, T., Morgenstern, J., Vecchione, B., Vaughan, J. W., Wallach, H., Daumé III, H., Crawford, K.: Datasheets for Datasets. \emph{Communications of the ACM} 64(12), 86--92, 2021.
\bibitem{howard1966} Howard, R. A.: Information Value Theory. \emph{IEEE Transactions on Systems Science and Cybernetics} 2(1), 22--26, 1966.
\bibitem{iso25012} ISO/IEC 25012:2008: \emph{Software engineering -- Software product Quality Requirements and Evaluation (SQuaRE) -- Data quality model}. International Organization for Standardization, 2008.
\bibitem{iso8000} ISO 8000-61:2016: \emph{Data quality -- Part 61: Data quality management: Process reference model}. International Organization for Standardization, 2016.
\bibitem{ji2023} Ji, Z., Lee, N., Frieske, R., Yu, T., Su, D., Xu, Y., Ishii, E., Bang, Y. J., Madotto, A., Fung, P.: Survey of Hallucination in Natural Language Generation. \emph{ACM Computing Surveys} 55(12), Article 248, 2023.
\bibitem{khatri2010} Khatri, V., Brown, C. V.: Designing Data Governance. \emph{Communications of the ACM} 53(1), 148--152, 2010.
\bibitem{lewis2020} Lewis, P., Perez, E., Piktus, A., Petroni, F., Karpukhin, V., Goyal, N., Küttler, H., Lewis, M., Yih, W., Rocktäschel, T., Riedel, S., Kiela, D.: Retrieval-Augmented Generation for Knowledge-Intensive NLP Tasks. In: \emph{Advances in Neural Information Processing Systems} 33, 9459--9474, 2020.
\bibitem{mitchell2019} Mitchell, M., Wu, S., Zaldivar, A., Barnes, P., Vasserman, L., Hutchinson, B., Spitzer, E., Raji, I. D., Gebru, T.: Model Cards for Model Reporting. In: \emph{Proceedings of the Conference on Fairness, Accountability, and Transparency (FAT* '19)}, 220--229, ACM, 2019.
\bibitem{moreau2011} Moreau, L., Clifford, B., Freire, J., Futrelle, J., Gil, Y., Groth, P., Kwasnikowska, N., Miles, S., Missier, P., Myers, J., Plale, B., Simmhan, Y., Stephan, E., Van den Bussche, J.: The Open Provenance Model core specification (v1.1). \emph{Future Generation Computer Systems} 27(6), 743--756, 2011.
\bibitem{nist2023} National Institute of Standards and Technology: \emph{Artificial Intelligence Risk Management Framework (AI RMF 1.0)}. NIST AI 100-1, Januar 2023. \url{https://doi.org/10.6028/NIST.AI.100-1}.
\bibitem{pipino2002} Pipino, L. L., Lee, Y. W., Wang, R. Y.: Data Quality Assessment. \emph{Communications of the ACM} 45(4), 211--218, 2002.
\bibitem{polyzotis2018} Polyzotis, N., Roy, S., Whang, S. E., Zinkevich, M.: Data Lifecycle Challenges in Production Machine Learning: A Survey. \emph{ACM SIGMOD Record} 47(2), 17--28, 2018.
\bibitem{prov2013} Lebo, T., Sahoo, S., McGuinness, D. (Hrsg.): \emph{PROV-O: The PROV Ontology}. W3C Recommendation, 30. April 2013. \url{https://www.w3.org/TR/prov-o/}.
\bibitem{raiffa1961} Raiffa, H., Schlaifer, R.: \emph{Applied Statistical Decision Theory}. Division of Research, Graduate School of Business Administration, Harvard University, 1961.
\bibitem{redman2016} Redman, T. C.: Bad Data Costs the U.S. \$3 Trillion Per Year. \emph{Harvard Business Review}, 22. September 2016.
\bibitem{sambasivan2021} Sambasivan, N., Kapania, S., Highfill, H., Akrong, D., Paritosh, P., Aroyo, L. M.: "Everyone wants to do the model work, not the data work": Data Cascades in High-Stakes AI. In: \emph{Proceedings of the CHI Conference on Human Factors in Computing Systems}, Article 39, ACM, 2021.
\bibitem{schelter2018} Schelter, S., Lange, D., Schmidt, P., Celikel, M., Biessmann, F., Grafberger, A.: Automating Large-Scale Data Quality Verification. \emph{Proceedings of the VLDB Endowment} 11(12), 1781--1794, 2018.
\bibitem{sculley2015} Sculley, D., Holt, G., Golovin, D., Davydov, E., Phillips, T., Ebner, D., Chaudhary, V., Young, M., Crespo, J.-F., Dennison, D.: Hidden Technical Debt in Machine Learning Systems. In: \emph{Advances in Neural Information Processing Systems} 28, 2503--2511, 2015.
\bibitem{wang1996} Wang, R. Y., Strong, D. M.: Beyond Accuracy: What Data Quality Means to Data Consumers. \emph{Journal of Management Information Systems} 12(4), 5--33, 1996.
\bibitem{whang2023} Whang, S. E., Roh, Y., Song, H., Lee, J.-G.: Data collection and quality challenges in deep learning: a data-centric AI perspective. \emph{The VLDB Journal} 32(4), 791--813, 2023.
\bibitem{wilkinson2016} Wilkinson, M. D., Dumontier, M., Aalbersberg, I. J., et al.: The FAIR Guiding Principles for scientific data management and stewardship. \emph{Scientific Data} 3, 160018, 2016.
\bibitem{zheng2023} Zheng, L., Chiang, W.-L., Sheng, Y., Zhuang, S., Wu, Z., Zhuang, Y., Lin, Z., Li, Z., Li, D., Xing, E. P., Zhang, H., Gonzalez, J. E., Stoica, I.: Judging LLM-as-a-Judge with MT-Bench and Chatbot Arena. In: \emph{Advances in Neural Information Processing Systems} 36 (Datasets and Benchmarks Track), 2023.
\end{thebibliography}

\end{document}
